
\documentclass[a4paper,12pt]{article}
\usepackage[margin=1cm]{geometry}
\usepackage{tikz}
\usepackage{qrcode}
\usepackage{multicol}
\usepackage{array}
\usepackage{adjustbox}
\pagestyle{empty}

% --- Bolinha centralizada ---
\newcommand{\marcacao}{%
    \vfill
    \tikz[baseline=(c.base)] \node (c) [circle, draw, minimum width=0.45cm, minimum height=0.45cm] {};
    \vfill
}

% --- Nova coluna centralizada ---
\newcolumntype{C}[1]{>{\centering\arraybackslash}m{#1}}

% --- Gabarito com 25 questões compactado ---
\newcommand{\Gabarito}{
    \renewcommand{\arraystretch}{1.2}
    \begin{center}
        \begin{tikzpicture}
            \node (table) [inner sep=0pt, anchor=north west] {
                \begin{tabular}{|C{0.6cm}|C{0.8cm}|C{0.8cm}|C{0.8cm}|C{0.8cm}|}
                    \hline
                    \adjustbox{max width=0.5cm}{\textbf{Nº}} & \textbf{A} & \textbf{B} & \textbf{C} & \textbf{D} \\
                    \hline
                    01 & \marcacao & \marcacao & \marcacao & \marcacao \\ \hline
02 & \marcacao & \marcacao & \marcacao & \marcacao \\ \hline
03 & \marcacao & \marcacao & \marcacao & \marcacao \\ \hline
04 & \marcacao & \marcacao & \marcacao & \marcacao \\ \hline
05 & \marcacao & \marcacao & \marcacao & \marcacao \\ \hline
06 & \marcacao & \marcacao & \marcacao & \marcacao \\ \hline
07 & \marcacao & \marcacao & \marcacao & \marcacao \\ \hline
08 & \marcacao & \marcacao & \marcacao & \marcacao \\ \hline
09 & \marcacao & \marcacao & \marcacao & \marcacao \\ \hline
10 & \marcacao & \marcacao & \marcacao & \marcacao \\ \hline
11 & \marcacao & \marcacao & \marcacao & \marcacao \\ \hline
12 & \marcacao & \marcacao & \marcacao & \marcacao \\ \hline
13 & \marcacao & \marcacao & \marcacao & \marcacao \\ \hline
14 & \marcacao & \marcacao & \marcacao & \marcacao \\ \hline
15 & \marcacao & \marcacao & \marcacao & \marcacao \\ \hline
16 & \marcacao & \marcacao & \marcacao & \marcacao \\ \hline
17 & \marcacao & \marcacao & \marcacao & \marcacao \\ \hline
18 & \marcacao & \marcacao & \marcacao & \marcacao \\ \hline
19 & \marcacao & \marcacao & \marcacao & \marcacao \\ \hline
20 & \marcacao & \marcacao & \marcacao & \marcacao \\ \hline
21 & \marcacao & \marcacao & \marcacao & \marcacao \\ \hline
22 & \marcacao & \marcacao & \marcacao & \marcacao \\ \hline
23 & \marcacao & \marcacao & \marcacao & \marcacao \\ \hline
24 & \marcacao & \marcacao & \marcacao & \marcacao \\ \hline
25 & \marcacao & \marcacao & \marcacao & \marcacao \\ \hline
                \end{tabular}
            };
            % --- Marcadores OMR ---
            \fill[black] ([shift={(-7mm,7mm)}]table.north west) rectangle ++(5mm,-5mm);
            \fill[black] ([shift={(7mm,7mm)}]table.north east) rectangle ++(-5mm,-5mm);
            \fill[black] ([shift={(-7mm,-7mm)}]table.south west) rectangle ++(5mm,5mm);
            \fill[black] ([shift={(7mm,-7mm)}]table.south east) rectangle ++(-5mm,5mm);
        \end{tikzpicture}
    \end{center}
    \renewcommand{\arraystretch}{1.0}
}

\begin{document}
% --- Bolinha centralizada ---
\newcommand{\marcacao}{%
    \vfill
    \tikz[baseline=(c.base)] \node (c) [circle, draw, minimum width=0.45cm, minimum height=0.45cm] {};
    \vfill
}

% --- Nova coluna centralizada ---
\newcolumntype{C}[1]{>{\centering\arraybackslash}m{#1}}

% --- Cabeçalho ---
\begin{minipage}[t]{0.7\textwidth}
    \textbf{Escola:} CEF04\_PLAN \\
    \textbf{Turma:} 2S - 9º ANO - Matutino \\
    \textbf{Aluno:} ÁGATA BEATRIZ DE SOUSA VIANA
\end{minipage}
\begin{minipage}[t]{0.29\textwidth}
    \raggedleft
    \qrcode[height=2.2cm]{ 462541-ÁGATA_BEATRIZ_DE_SOUSA_VIANA-2S }
\end{minipage}

\vspace{0.5cm}
\hrule
\vspace{0.5cm}

% --- Gabarito com 25 questões compactado ---
\newcommand{\Gabarito}{
    \renewcommand{\arraystretch}{1.2} % altura das linhas compacta
    \begin{center}
        \begin{tikzpicture}
            \node (table) [inner sep=0pt, anchor=north west] {
                \begin{tabular}{|C{0.6cm}|C{0.8cm}|C{0.8cm}|C{0.8cm}|C{0.8cm}|}
                    \hline
                    \adjustbox{max width=0.5cm}{\textbf{Nº}} & \textbf{A} & \textbf{B} & \textbf{C} & \textbf{D} \\
                    \hline
                    01 & \marcacao & \marcacao & \marcacao & \marcacao \\ \hline
                    02 & \marcacao & \marcacao & \marcacao & \marcacao \\ \hline
                    03 & \marcacao & \marcacao & \marcacao & \marcacao \\ \hline
                    04 & \marcacao & \marcacao & \marcacao & \marcacao \\ \hline
                    05 & \marcacao & \marcacao & \marcacao & \marcacao \\ \hline
                    06 & \marcacao & \marcacao & \marcacao & \marcacao \\ \hline
                    07 & \marcacao & \marcacao & \marcacao & \marcacao \\ \hline
                    08 & \marcacao & \marcacao & \marcacao & \marcacao \\ \hline
                    09 & \marcacao & \marcacao & \marcacao & \marcacao \\ \hline
                    10 & \marcacao & \marcacao & \marcacao & \marcacao \\ \hline
                    11 & \marcacao & \marcacao & \marcacao & \marcacao \\ \hline
                    12 & \marcacao & \marcacao & \marcacao & \marcacao \\ \hline
                    13 & \marcacao & \marcacao & \marcacao & \marcacao \\ \hline
                    14 & \marcacao & \marcacao & \marcacao & \marcacao \\ \hline
                    15 & \marcacao & \marcacao & \marcacao & \marcacao \\ \hline
                    16 & \marcacao & \marcacao & \marcacao & \marcacao \\ \hline
                    17 & \marcacao & \marcacao & \marcacao & \marcacao \\ \hline
                    18 & \marcacao & \marcacao & \marcacao & \marcacao \\ \hline
                    19 & \marcacao & \marcacao & \marcacao & \marcacao \\ \hline
                    20 & \marcacao & \marcacao & \marcacao & \marcacao \\ \hline
                    21 & \marcacao & \marcacao & \marcacao & \marcacao \\ \hline
                    22 & \marcacao & \marcacao & \marcacao & \marcacao \\ \hline
                    23 & \marcacao & \marcacao & \marcacao & \marcacao \\ \hline
                    24 & \marcacao & \marcacao & \marcacao & \marcacao \\ \hline
                    25 & \marcacao & \marcacao & \marcacao & \marcacao \\ \hline
                \end{tabular}
            };
            
            % --- Marcadores OMR ---
            \fill[black] ([shift={(-7mm,7mm)}]table.north west) rectangle ++(5mm,-5mm);
            \fill[black] ([shift={(7mm,7mm)}]table.north east) rectangle ++(-5mm,-5mm);
            \fill[black] ([shift={(-7mm,-7mm)}]table.south west) rectangle ++(5mm,5mm);
            \fill[black] ([shift={(7mm,-7mm)}]table.south east) rectangle ++(-5mm,5mm);
        \end{tikzpicture}
    \end{center}
    \renewcommand{\arraystretch}{1.0}
}

% Chamada do gabarito
\Gabarito

\newpage
% --- Bolinha centralizada ---
\newcommand{\marcacao}{%
    \vfill
    \tikz[baseline=(c.base)] \node (c) [circle, draw, minimum width=0.45cm, minimum height=0.45cm] {};
    \vfill
}

% --- Nova coluna centralizada ---
\newcolumntype{C}[1]{>{\centering\arraybackslash}m{#1}}

% --- Cabeçalho ---
\begin{minipage}[t]{0.7\textwidth}
    \textbf{Escola:} CEF04\_PLAN \\
    \textbf{Turma:} 2S - 9º ANO - Matutino \\
    \textbf{Aluno:} ANDRE PEREIRA DOS SANTOS
\end{minipage}
\begin{minipage}[t]{0.29\textwidth}
    \raggedleft
    \qrcode[height=2.2cm]{ 326187-ANDRE_PEREIRA_DOS_SANTOS-2S }
\end{minipage}

\vspace{0.5cm}
\hrule
\vspace{0.5cm}

% --- Gabarito com 25 questões compactado ---
\newcommand{\Gabarito}{
    \renewcommand{\arraystretch}{1.2} % altura das linhas compacta
    \begin{center}
        \begin{tikzpicture}
            \node (table) [inner sep=0pt, anchor=north west] {
                \begin{tabular}{|C{0.6cm}|C{0.8cm}|C{0.8cm}|C{0.8cm}|C{0.8cm}|}
                    \hline
                    \adjustbox{max width=0.5cm}{\textbf{Nº}} & \textbf{A} & \textbf{B} & \textbf{C} & \textbf{D} \\
                    \hline
                    01 & \marcacao & \marcacao & \marcacao & \marcacao \\ \hline
                    02 & \marcacao & \marcacao & \marcacao & \marcacao \\ \hline
                    03 & \marcacao & \marcacao & \marcacao & \marcacao \\ \hline
                    04 & \marcacao & \marcacao & \marcacao & \marcacao \\ \hline
                    05 & \marcacao & \marcacao & \marcacao & \marcacao \\ \hline
                    06 & \marcacao & \marcacao & \marcacao & \marcacao \\ \hline
                    07 & \marcacao & \marcacao & \marcacao & \marcacao \\ \hline
                    08 & \marcacao & \marcacao & \marcacao & \marcacao \\ \hline
                    09 & \marcacao & \marcacao & \marcacao & \marcacao \\ \hline
                    10 & \marcacao & \marcacao & \marcacao & \marcacao \\ \hline
                    11 & \marcacao & \marcacao & \marcacao & \marcacao \\ \hline
                    12 & \marcacao & \marcacao & \marcacao & \marcacao \\ \hline
                    13 & \marcacao & \marcacao & \marcacao & \marcacao \\ \hline
                    14 & \marcacao & \marcacao & \marcacao & \marcacao \\ \hline
                    15 & \marcacao & \marcacao & \marcacao & \marcacao \\ \hline
                    16 & \marcacao & \marcacao & \marcacao & \marcacao \\ \hline
                    17 & \marcacao & \marcacao & \marcacao & \marcacao \\ \hline
                    18 & \marcacao & \marcacao & \marcacao & \marcacao \\ \hline
                    19 & \marcacao & \marcacao & \marcacao & \marcacao \\ \hline
                    20 & \marcacao & \marcacao & \marcacao & \marcacao \\ \hline
                    21 & \marcacao & \marcacao & \marcacao & \marcacao \\ \hline
                    22 & \marcacao & \marcacao & \marcacao & \marcacao \\ \hline
                    23 & \marcacao & \marcacao & \marcacao & \marcacao \\ \hline
                    24 & \marcacao & \marcacao & \marcacao & \marcacao \\ \hline
                    25 & \marcacao & \marcacao & \marcacao & \marcacao \\ \hline
                \end{tabular}
            };
            
            % --- Marcadores OMR ---
            \fill[black] ([shift={(-7mm,7mm)}]table.north west) rectangle ++(5mm,-5mm);
            \fill[black] ([shift={(7mm,7mm)}]table.north east) rectangle ++(-5mm,-5mm);
            \fill[black] ([shift={(-7mm,-7mm)}]table.south west) rectangle ++(5mm,5mm);
            \fill[black] ([shift={(7mm,-7mm)}]table.south east) rectangle ++(-5mm,5mm);
        \end{tikzpicture}
    \end{center}
    \renewcommand{\arraystretch}{1.0}
}

% Chamada do gabarito
\Gabarito

\newpage
% --- Bolinha centralizada ---
\newcommand{\marcacao}{%
    \vfill
    \tikz[baseline=(c.base)] \node (c) [circle, draw, minimum width=0.45cm, minimum height=0.45cm] {};
    \vfill
}

% --- Nova coluna centralizada ---
\newcolumntype{C}[1]{>{\centering\arraybackslash}m{#1}}

% --- Cabeçalho ---
\begin{minipage}[t]{0.7\textwidth}
    \textbf{Escola:} CEF04\_PLAN \\
    \textbf{Turma:} 2S - 9º ANO - Matutino \\
    \textbf{Aluno:} ANDRESON MENDONÇA DE ANDRADE
\end{minipage}
\begin{minipage}[t]{0.29\textwidth}
    \raggedleft
    \qrcode[height=2.2cm]{ 477427-ANDRESON_MENDONÇA_DE_ANDRADE-2S }
\end{minipage}

\vspace{0.5cm}
\hrule
\vspace{0.5cm}

% --- Gabarito com 25 questões compactado ---
\newcommand{\Gabarito}{
    \renewcommand{\arraystretch}{1.2} % altura das linhas compacta
    \begin{center}
        \begin{tikzpicture}
            \node (table) [inner sep=0pt, anchor=north west] {
                \begin{tabular}{|C{0.6cm}|C{0.8cm}|C{0.8cm}|C{0.8cm}|C{0.8cm}|}
                    \hline
                    \adjustbox{max width=0.5cm}{\textbf{Nº}} & \textbf{A} & \textbf{B} & \textbf{C} & \textbf{D} \\
                    \hline
                    01 & \marcacao & \marcacao & \marcacao & \marcacao \\ \hline
                    02 & \marcacao & \marcacao & \marcacao & \marcacao \\ \hline
                    03 & \marcacao & \marcacao & \marcacao & \marcacao \\ \hline
                    04 & \marcacao & \marcacao & \marcacao & \marcacao \\ \hline
                    05 & \marcacao & \marcacao & \marcacao & \marcacao \\ \hline
                    06 & \marcacao & \marcacao & \marcacao & \marcacao \\ \hline
                    07 & \marcacao & \marcacao & \marcacao & \marcacao \\ \hline
                    08 & \marcacao & \marcacao & \marcacao & \marcacao \\ \hline
                    09 & \marcacao & \marcacao & \marcacao & \marcacao \\ \hline
                    10 & \marcacao & \marcacao & \marcacao & \marcacao \\ \hline
                    11 & \marcacao & \marcacao & \marcacao & \marcacao \\ \hline
                    12 & \marcacao & \marcacao & \marcacao & \marcacao \\ \hline
                    13 & \marcacao & \marcacao & \marcacao & \marcacao \\ \hline
                    14 & \marcacao & \marcacao & \marcacao & \marcacao \\ \hline
                    15 & \marcacao & \marcacao & \marcacao & \marcacao \\ \hline
                    16 & \marcacao & \marcacao & \marcacao & \marcacao \\ \hline
                    17 & \marcacao & \marcacao & \marcacao & \marcacao \\ \hline
                    18 & \marcacao & \marcacao & \marcacao & \marcacao \\ \hline
                    19 & \marcacao & \marcacao & \marcacao & \marcacao \\ \hline
                    20 & \marcacao & \marcacao & \marcacao & \marcacao \\ \hline
                    21 & \marcacao & \marcacao & \marcacao & \marcacao \\ \hline
                    22 & \marcacao & \marcacao & \marcacao & \marcacao \\ \hline
                    23 & \marcacao & \marcacao & \marcacao & \marcacao \\ \hline
                    24 & \marcacao & \marcacao & \marcacao & \marcacao \\ \hline
                    25 & \marcacao & \marcacao & \marcacao & \marcacao \\ \hline
                \end{tabular}
            };
            
            % --- Marcadores OMR ---
            \fill[black] ([shift={(-7mm,7mm)}]table.north west) rectangle ++(5mm,-5mm);
            \fill[black] ([shift={(7mm,7mm)}]table.north east) rectangle ++(-5mm,-5mm);
            \fill[black] ([shift={(-7mm,-7mm)}]table.south west) rectangle ++(5mm,5mm);
            \fill[black] ([shift={(7mm,-7mm)}]table.south east) rectangle ++(-5mm,5mm);
        \end{tikzpicture}
    \end{center}
    \renewcommand{\arraystretch}{1.0}
}

% Chamada do gabarito
\Gabarito

\newpage
% --- Bolinha centralizada ---
\newcommand{\marcacao}{%
    \vfill
    \tikz[baseline=(c.base)] \node (c) [circle, draw, minimum width=0.45cm, minimum height=0.45cm] {};
    \vfill
}

% --- Nova coluna centralizada ---
\newcolumntype{C}[1]{>{\centering\arraybackslash}m{#1}}

% --- Cabeçalho ---
\begin{minipage}[t]{0.7\textwidth}
    \textbf{Escola:} CEF04\_PLAN \\
    \textbf{Turma:} 2S - 9º ANO - Matutino \\
    \textbf{Aluno:} ARTHUR MIGUEL DIAS SILVA
\end{minipage}
\begin{minipage}[t]{0.29\textwidth}
    \raggedleft
    \qrcode[height=2.2cm]{ 425199-ARTHUR_MIGUEL_DIAS_SILVA-2S }
\end{minipage}

\vspace{0.5cm}
\hrule
\vspace{0.5cm}

% --- Gabarito com 25 questões compactado ---
\newcommand{\Gabarito}{
    \renewcommand{\arraystretch}{1.2} % altura das linhas compacta
    \begin{center}
        \begin{tikzpicture}
            \node (table) [inner sep=0pt, anchor=north west] {
                \begin{tabular}{|C{0.6cm}|C{0.8cm}|C{0.8cm}|C{0.8cm}|C{0.8cm}|}
                    \hline
                    \adjustbox{max width=0.5cm}{\textbf{Nº}} & \textbf{A} & \textbf{B} & \textbf{C} & \textbf{D} \\
                    \hline
                    01 & \marcacao & \marcacao & \marcacao & \marcacao \\ \hline
                    02 & \marcacao & \marcacao & \marcacao & \marcacao \\ \hline
                    03 & \marcacao & \marcacao & \marcacao & \marcacao \\ \hline
                    04 & \marcacao & \marcacao & \marcacao & \marcacao \\ \hline
                    05 & \marcacao & \marcacao & \marcacao & \marcacao \\ \hline
                    06 & \marcacao & \marcacao & \marcacao & \marcacao \\ \hline
                    07 & \marcacao & \marcacao & \marcacao & \marcacao \\ \hline
                    08 & \marcacao & \marcacao & \marcacao & \marcacao \\ \hline
                    09 & \marcacao & \marcacao & \marcacao & \marcacao \\ \hline
                    10 & \marcacao & \marcacao & \marcacao & \marcacao \\ \hline
                    11 & \marcacao & \marcacao & \marcacao & \marcacao \\ \hline
                    12 & \marcacao & \marcacao & \marcacao & \marcacao \\ \hline
                    13 & \marcacao & \marcacao & \marcacao & \marcacao \\ \hline
                    14 & \marcacao & \marcacao & \marcacao & \marcacao \\ \hline
                    15 & \marcacao & \marcacao & \marcacao & \marcacao \\ \hline
                    16 & \marcacao & \marcacao & \marcacao & \marcacao \\ \hline
                    17 & \marcacao & \marcacao & \marcacao & \marcacao \\ \hline
                    18 & \marcacao & \marcacao & \marcacao & \marcacao \\ \hline
                    19 & \marcacao & \marcacao & \marcacao & \marcacao \\ \hline
                    20 & \marcacao & \marcacao & \marcacao & \marcacao \\ \hline
                    21 & \marcacao & \marcacao & \marcacao & \marcacao \\ \hline
                    22 & \marcacao & \marcacao & \marcacao & \marcacao \\ \hline
                    23 & \marcacao & \marcacao & \marcacao & \marcacao \\ \hline
                    24 & \marcacao & \marcacao & \marcacao & \marcacao \\ \hline
                    25 & \marcacao & \marcacao & \marcacao & \marcacao \\ \hline
                \end{tabular}
            };
            
            % --- Marcadores OMR ---
            \fill[black] ([shift={(-7mm,7mm)}]table.north west) rectangle ++(5mm,-5mm);
            \fill[black] ([shift={(7mm,7mm)}]table.north east) rectangle ++(-5mm,-5mm);
            \fill[black] ([shift={(-7mm,-7mm)}]table.south west) rectangle ++(5mm,5mm);
            \fill[black] ([shift={(7mm,-7mm)}]table.south east) rectangle ++(-5mm,5mm);
        \end{tikzpicture}
    \end{center}
    \renewcommand{\arraystretch}{1.0}
}

% Chamada do gabarito
\Gabarito

\newpage
% --- Bolinha centralizada ---
\newcommand{\marcacao}{%
    \vfill
    \tikz[baseline=(c.base)] \node (c) [circle, draw, minimum width=0.45cm, minimum height=0.45cm] {};
    \vfill
}

% --- Nova coluna centralizada ---
\newcolumntype{C}[1]{>{\centering\arraybackslash}m{#1}}

% --- Cabeçalho ---
\begin{minipage}[t]{0.7\textwidth}
    \textbf{Escola:} CEF04\_PLAN \\
    \textbf{Turma:} 2S - 9º ANO - Matutino \\
    \textbf{Aluno:} BRUNA PIRES MAIA ALVES
\end{minipage}
\begin{minipage}[t]{0.29\textwidth}
    \raggedleft
    \qrcode[height=2.2cm]{ 99570-BRUNA_PIRES_MAIA_ALVES-2S }
\end{minipage}

\vspace{0.5cm}
\hrule
\vspace{0.5cm}

% --- Gabarito com 25 questões compactado ---
\newcommand{\Gabarito}{
    \renewcommand{\arraystretch}{1.2} % altura das linhas compacta
    \begin{center}
        \begin{tikzpicture}
            \node (table) [inner sep=0pt, anchor=north west] {
                \begin{tabular}{|C{0.6cm}|C{0.8cm}|C{0.8cm}|C{0.8cm}|C{0.8cm}|}
                    \hline
                    \adjustbox{max width=0.5cm}{\textbf{Nº}} & \textbf{A} & \textbf{B} & \textbf{C} & \textbf{D} \\
                    \hline
                    01 & \marcacao & \marcacao & \marcacao & \marcacao \\ \hline
                    02 & \marcacao & \marcacao & \marcacao & \marcacao \\ \hline
                    03 & \marcacao & \marcacao & \marcacao & \marcacao \\ \hline
                    04 & \marcacao & \marcacao & \marcacao & \marcacao \\ \hline
                    05 & \marcacao & \marcacao & \marcacao & \marcacao \\ \hline
                    06 & \marcacao & \marcacao & \marcacao & \marcacao \\ \hline
                    07 & \marcacao & \marcacao & \marcacao & \marcacao \\ \hline
                    08 & \marcacao & \marcacao & \marcacao & \marcacao \\ \hline
                    09 & \marcacao & \marcacao & \marcacao & \marcacao \\ \hline
                    10 & \marcacao & \marcacao & \marcacao & \marcacao \\ \hline
                    11 & \marcacao & \marcacao & \marcacao & \marcacao \\ \hline
                    12 & \marcacao & \marcacao & \marcacao & \marcacao \\ \hline
                    13 & \marcacao & \marcacao & \marcacao & \marcacao \\ \hline
                    14 & \marcacao & \marcacao & \marcacao & \marcacao \\ \hline
                    15 & \marcacao & \marcacao & \marcacao & \marcacao \\ \hline
                    16 & \marcacao & \marcacao & \marcacao & \marcacao \\ \hline
                    17 & \marcacao & \marcacao & \marcacao & \marcacao \\ \hline
                    18 & \marcacao & \marcacao & \marcacao & \marcacao \\ \hline
                    19 & \marcacao & \marcacao & \marcacao & \marcacao \\ \hline
                    20 & \marcacao & \marcacao & \marcacao & \marcacao \\ \hline
                    21 & \marcacao & \marcacao & \marcacao & \marcacao \\ \hline
                    22 & \marcacao & \marcacao & \marcacao & \marcacao \\ \hline
                    23 & \marcacao & \marcacao & \marcacao & \marcacao \\ \hline
                    24 & \marcacao & \marcacao & \marcacao & \marcacao \\ \hline
                    25 & \marcacao & \marcacao & \marcacao & \marcacao \\ \hline
                \end{tabular}
            };
            
            % --- Marcadores OMR ---
            \fill[black] ([shift={(-7mm,7mm)}]table.north west) rectangle ++(5mm,-5mm);
            \fill[black] ([shift={(7mm,7mm)}]table.north east) rectangle ++(-5mm,-5mm);
            \fill[black] ([shift={(-7mm,-7mm)}]table.south west) rectangle ++(5mm,5mm);
            \fill[black] ([shift={(7mm,-7mm)}]table.south east) rectangle ++(-5mm,5mm);
        \end{tikzpicture}
    \end{center}
    \renewcommand{\arraystretch}{1.0}
}

% Chamada do gabarito
\Gabarito

\newpage
% --- Bolinha centralizada ---
\newcommand{\marcacao}{%
    \vfill
    \tikz[baseline=(c.base)] \node (c) [circle, draw, minimum width=0.45cm, minimum height=0.45cm] {};
    \vfill
}

% --- Nova coluna centralizada ---
\newcolumntype{C}[1]{>{\centering\arraybackslash}m{#1}}

% --- Cabeçalho ---
\begin{minipage}[t]{0.7\textwidth}
    \textbf{Escola:} CEF04\_PLAN \\
    \textbf{Turma:} 2S - 9º ANO - Matutino \\
    \textbf{Aluno:} DANIEL GOMES FERNANDES
\end{minipage}
\begin{minipage}[t]{0.29\textwidth}
    \raggedleft
    \qrcode[height=2.2cm]{ 492984-DANIEL_GOMES_FERNANDES-2S }
\end{minipage}

\vspace{0.5cm}
\hrule
\vspace{0.5cm}

% --- Gabarito com 25 questões compactado ---
\newcommand{\Gabarito}{
    \renewcommand{\arraystretch}{1.2} % altura das linhas compacta
    \begin{center}
        \begin{tikzpicture}
            \node (table) [inner sep=0pt, anchor=north west] {
                \begin{tabular}{|C{0.6cm}|C{0.8cm}|C{0.8cm}|C{0.8cm}|C{0.8cm}|}
                    \hline
                    \adjustbox{max width=0.5cm}{\textbf{Nº}} & \textbf{A} & \textbf{B} & \textbf{C} & \textbf{D} \\
                    \hline
                    01 & \marcacao & \marcacao & \marcacao & \marcacao \\ \hline
                    02 & \marcacao & \marcacao & \marcacao & \marcacao \\ \hline
                    03 & \marcacao & \marcacao & \marcacao & \marcacao \\ \hline
                    04 & \marcacao & \marcacao & \marcacao & \marcacao \\ \hline
                    05 & \marcacao & \marcacao & \marcacao & \marcacao \\ \hline
                    06 & \marcacao & \marcacao & \marcacao & \marcacao \\ \hline
                    07 & \marcacao & \marcacao & \marcacao & \marcacao \\ \hline
                    08 & \marcacao & \marcacao & \marcacao & \marcacao \\ \hline
                    09 & \marcacao & \marcacao & \marcacao & \marcacao \\ \hline
                    10 & \marcacao & \marcacao & \marcacao & \marcacao \\ \hline
                    11 & \marcacao & \marcacao & \marcacao & \marcacao \\ \hline
                    12 & \marcacao & \marcacao & \marcacao & \marcacao \\ \hline
                    13 & \marcacao & \marcacao & \marcacao & \marcacao \\ \hline
                    14 & \marcacao & \marcacao & \marcacao & \marcacao \\ \hline
                    15 & \marcacao & \marcacao & \marcacao & \marcacao \\ \hline
                    16 & \marcacao & \marcacao & \marcacao & \marcacao \\ \hline
                    17 & \marcacao & \marcacao & \marcacao & \marcacao \\ \hline
                    18 & \marcacao & \marcacao & \marcacao & \marcacao \\ \hline
                    19 & \marcacao & \marcacao & \marcacao & \marcacao \\ \hline
                    20 & \marcacao & \marcacao & \marcacao & \marcacao \\ \hline
                    21 & \marcacao & \marcacao & \marcacao & \marcacao \\ \hline
                    22 & \marcacao & \marcacao & \marcacao & \marcacao \\ \hline
                    23 & \marcacao & \marcacao & \marcacao & \marcacao \\ \hline
                    24 & \marcacao & \marcacao & \marcacao & \marcacao \\ \hline
                    25 & \marcacao & \marcacao & \marcacao & \marcacao \\ \hline
                \end{tabular}
            };
            
            % --- Marcadores OMR ---
            \fill[black] ([shift={(-7mm,7mm)}]table.north west) rectangle ++(5mm,-5mm);
            \fill[black] ([shift={(7mm,7mm)}]table.north east) rectangle ++(-5mm,-5mm);
            \fill[black] ([shift={(-7mm,-7mm)}]table.south west) rectangle ++(5mm,5mm);
            \fill[black] ([shift={(7mm,-7mm)}]table.south east) rectangle ++(-5mm,5mm);
        \end{tikzpicture}
    \end{center}
    \renewcommand{\arraystretch}{1.0}
}

% Chamada do gabarito
\Gabarito

\newpage
% --- Bolinha centralizada ---
\newcommand{\marcacao}{%
    \vfill
    \tikz[baseline=(c.base)] \node (c) [circle, draw, minimum width=0.45cm, minimum height=0.45cm] {};
    \vfill
}

% --- Nova coluna centralizada ---
\newcolumntype{C}[1]{>{\centering\arraybackslash}m{#1}}

% --- Cabeçalho ---
\begin{minipage}[t]{0.7\textwidth}
    \textbf{Escola:} CEF04\_PLAN \\
    \textbf{Turma:} 2S - 9º ANO - Matutino \\
    \textbf{Aluno:} EDUARDO PHILIPE LIMA RIBEIRO
\end{minipage}
\begin{minipage}[t]{0.29\textwidth}
    \raggedleft
    \qrcode[height=2.2cm]{ 1211137-EDUARDO_PHILIPE_LIMA_RIBEIRO-2S }
\end{minipage}

\vspace{0.5cm}
\hrule
\vspace{0.5cm}

% --- Gabarito com 25 questões compactado ---
\newcommand{\Gabarito}{
    \renewcommand{\arraystretch}{1.2} % altura das linhas compacta
    \begin{center}
        \begin{tikzpicture}
            \node (table) [inner sep=0pt, anchor=north west] {
                \begin{tabular}{|C{0.6cm}|C{0.8cm}|C{0.8cm}|C{0.8cm}|C{0.8cm}|}
                    \hline
                    \adjustbox{max width=0.5cm}{\textbf{Nº}} & \textbf{A} & \textbf{B} & \textbf{C} & \textbf{D} \\
                    \hline
                    01 & \marcacao & \marcacao & \marcacao & \marcacao \\ \hline
                    02 & \marcacao & \marcacao & \marcacao & \marcacao \\ \hline
                    03 & \marcacao & \marcacao & \marcacao & \marcacao \\ \hline
                    04 & \marcacao & \marcacao & \marcacao & \marcacao \\ \hline
                    05 & \marcacao & \marcacao & \marcacao & \marcacao \\ \hline
                    06 & \marcacao & \marcacao & \marcacao & \marcacao \\ \hline
                    07 & \marcacao & \marcacao & \marcacao & \marcacao \\ \hline
                    08 & \marcacao & \marcacao & \marcacao & \marcacao \\ \hline
                    09 & \marcacao & \marcacao & \marcacao & \marcacao \\ \hline
                    10 & \marcacao & \marcacao & \marcacao & \marcacao \\ \hline
                    11 & \marcacao & \marcacao & \marcacao & \marcacao \\ \hline
                    12 & \marcacao & \marcacao & \marcacao & \marcacao \\ \hline
                    13 & \marcacao & \marcacao & \marcacao & \marcacao \\ \hline
                    14 & \marcacao & \marcacao & \marcacao & \marcacao \\ \hline
                    15 & \marcacao & \marcacao & \marcacao & \marcacao \\ \hline
                    16 & \marcacao & \marcacao & \marcacao & \marcacao \\ \hline
                    17 & \marcacao & \marcacao & \marcacao & \marcacao \\ \hline
                    18 & \marcacao & \marcacao & \marcacao & \marcacao \\ \hline
                    19 & \marcacao & \marcacao & \marcacao & \marcacao \\ \hline
                    20 & \marcacao & \marcacao & \marcacao & \marcacao \\ \hline
                    21 & \marcacao & \marcacao & \marcacao & \marcacao \\ \hline
                    22 & \marcacao & \marcacao & \marcacao & \marcacao \\ \hline
                    23 & \marcacao & \marcacao & \marcacao & \marcacao \\ \hline
                    24 & \marcacao & \marcacao & \marcacao & \marcacao \\ \hline
                    25 & \marcacao & \marcacao & \marcacao & \marcacao \\ \hline
                \end{tabular}
            };
            
            % --- Marcadores OMR ---
            \fill[black] ([shift={(-7mm,7mm)}]table.north west) rectangle ++(5mm,-5mm);
            \fill[black] ([shift={(7mm,7mm)}]table.north east) rectangle ++(-5mm,-5mm);
            \fill[black] ([shift={(-7mm,-7mm)}]table.south west) rectangle ++(5mm,5mm);
            \fill[black] ([shift={(7mm,-7mm)}]table.south east) rectangle ++(-5mm,5mm);
        \end{tikzpicture}
    \end{center}
    \renewcommand{\arraystretch}{1.0}
}

% Chamada do gabarito
\Gabarito

\newpage
% --- Bolinha centralizada ---
\newcommand{\marcacao}{%
    \vfill
    \tikz[baseline=(c.base)] \node (c) [circle, draw, minimum width=0.45cm, minimum height=0.45cm] {};
    \vfill
}

% --- Nova coluna centralizada ---
\newcolumntype{C}[1]{>{\centering\arraybackslash}m{#1}}

% --- Cabeçalho ---
\begin{minipage}[t]{0.7\textwidth}
    \textbf{Escola:} CEF04\_PLAN \\
    \textbf{Turma:} 2S - 9º ANO - Matutino \\
    \textbf{Aluno:} EMILLY VICTORIA LIMA DE SOUZA
\end{minipage}
\begin{minipage}[t]{0.29\textwidth}
    \raggedleft
    \qrcode[height=2.2cm]{ 518606-EMILLY_VICTORIA_LIMA_DE_SOUZA-2S }
\end{minipage}

\vspace{0.5cm}
\hrule
\vspace{0.5cm}

% --- Gabarito com 25 questões compactado ---
\newcommand{\Gabarito}{
    \renewcommand{\arraystretch}{1.2} % altura das linhas compacta
    \begin{center}
        \begin{tikzpicture}
            \node (table) [inner sep=0pt, anchor=north west] {
                \begin{tabular}{|C{0.6cm}|C{0.8cm}|C{0.8cm}|C{0.8cm}|C{0.8cm}|}
                    \hline
                    \adjustbox{max width=0.5cm}{\textbf{Nº}} & \textbf{A} & \textbf{B} & \textbf{C} & \textbf{D} \\
                    \hline
                    01 & \marcacao & \marcacao & \marcacao & \marcacao \\ \hline
                    02 & \marcacao & \marcacao & \marcacao & \marcacao \\ \hline
                    03 & \marcacao & \marcacao & \marcacao & \marcacao \\ \hline
                    04 & \marcacao & \marcacao & \marcacao & \marcacao \\ \hline
                    05 & \marcacao & \marcacao & \marcacao & \marcacao \\ \hline
                    06 & \marcacao & \marcacao & \marcacao & \marcacao \\ \hline
                    07 & \marcacao & \marcacao & \marcacao & \marcacao \\ \hline
                    08 & \marcacao & \marcacao & \marcacao & \marcacao \\ \hline
                    09 & \marcacao & \marcacao & \marcacao & \marcacao \\ \hline
                    10 & \marcacao & \marcacao & \marcacao & \marcacao \\ \hline
                    11 & \marcacao & \marcacao & \marcacao & \marcacao \\ \hline
                    12 & \marcacao & \marcacao & \marcacao & \marcacao \\ \hline
                    13 & \marcacao & \marcacao & \marcacao & \marcacao \\ \hline
                    14 & \marcacao & \marcacao & \marcacao & \marcacao \\ \hline
                    15 & \marcacao & \marcacao & \marcacao & \marcacao \\ \hline
                    16 & \marcacao & \marcacao & \marcacao & \marcacao \\ \hline
                    17 & \marcacao & \marcacao & \marcacao & \marcacao \\ \hline
                    18 & \marcacao & \marcacao & \marcacao & \marcacao \\ \hline
                    19 & \marcacao & \marcacao & \marcacao & \marcacao \\ \hline
                    20 & \marcacao & \marcacao & \marcacao & \marcacao \\ \hline
                    21 & \marcacao & \marcacao & \marcacao & \marcacao \\ \hline
                    22 & \marcacao & \marcacao & \marcacao & \marcacao \\ \hline
                    23 & \marcacao & \marcacao & \marcacao & \marcacao \\ \hline
                    24 & \marcacao & \marcacao & \marcacao & \marcacao \\ \hline
                    25 & \marcacao & \marcacao & \marcacao & \marcacao \\ \hline
                \end{tabular}
            };
            
            % --- Marcadores OMR ---
            \fill[black] ([shift={(-7mm,7mm)}]table.north west) rectangle ++(5mm,-5mm);
            \fill[black] ([shift={(7mm,7mm)}]table.north east) rectangle ++(-5mm,-5mm);
            \fill[black] ([shift={(-7mm,-7mm)}]table.south west) rectangle ++(5mm,5mm);
            \fill[black] ([shift={(7mm,-7mm)}]table.south east) rectangle ++(-5mm,5mm);
        \end{tikzpicture}
    \end{center}
    \renewcommand{\arraystretch}{1.0}
}

% Chamada do gabarito
\Gabarito

\newpage
% --- Bolinha centralizada ---
\newcommand{\marcacao}{%
    \vfill
    \tikz[baseline=(c.base)] \node (c) [circle, draw, minimum width=0.45cm, minimum height=0.45cm] {};
    \vfill
}

% --- Nova coluna centralizada ---
\newcolumntype{C}[1]{>{\centering\arraybackslash}m{#1}}

% --- Cabeçalho ---
\begin{minipage}[t]{0.7\textwidth}
    \textbf{Escola:} CEF04\_PLAN \\
    \textbf{Turma:} 2S - 9º ANO - Matutino \\
    \textbf{Aluno:} GUILHERME FERREIRA PONTE
\end{minipage}
\begin{minipage}[t]{0.29\textwidth}
    \raggedleft
    \qrcode[height=2.2cm]{ 493228-GUILHERME_FERREIRA_PONTE-2S }
\end{minipage}

\vspace{0.5cm}
\hrule
\vspace{0.5cm}

% --- Gabarito com 25 questões compactado ---
\newcommand{\Gabarito}{
    \renewcommand{\arraystretch}{1.2} % altura das linhas compacta
    \begin{center}
        \begin{tikzpicture}
            \node (table) [inner sep=0pt, anchor=north west] {
                \begin{tabular}{|C{0.6cm}|C{0.8cm}|C{0.8cm}|C{0.8cm}|C{0.8cm}|}
                    \hline
                    \adjustbox{max width=0.5cm}{\textbf{Nº}} & \textbf{A} & \textbf{B} & \textbf{C} & \textbf{D} \\
                    \hline
                    01 & \marcacao & \marcacao & \marcacao & \marcacao \\ \hline
                    02 & \marcacao & \marcacao & \marcacao & \marcacao \\ \hline
                    03 & \marcacao & \marcacao & \marcacao & \marcacao \\ \hline
                    04 & \marcacao & \marcacao & \marcacao & \marcacao \\ \hline
                    05 & \marcacao & \marcacao & \marcacao & \marcacao \\ \hline
                    06 & \marcacao & \marcacao & \marcacao & \marcacao \\ \hline
                    07 & \marcacao & \marcacao & \marcacao & \marcacao \\ \hline
                    08 & \marcacao & \marcacao & \marcacao & \marcacao \\ \hline
                    09 & \marcacao & \marcacao & \marcacao & \marcacao \\ \hline
                    10 & \marcacao & \marcacao & \marcacao & \marcacao \\ \hline
                    11 & \marcacao & \marcacao & \marcacao & \marcacao \\ \hline
                    12 & \marcacao & \marcacao & \marcacao & \marcacao \\ \hline
                    13 & \marcacao & \marcacao & \marcacao & \marcacao \\ \hline
                    14 & \marcacao & \marcacao & \marcacao & \marcacao \\ \hline
                    15 & \marcacao & \marcacao & \marcacao & \marcacao \\ \hline
                    16 & \marcacao & \marcacao & \marcacao & \marcacao \\ \hline
                    17 & \marcacao & \marcacao & \marcacao & \marcacao \\ \hline
                    18 & \marcacao & \marcacao & \marcacao & \marcacao \\ \hline
                    19 & \marcacao & \marcacao & \marcacao & \marcacao \\ \hline
                    20 & \marcacao & \marcacao & \marcacao & \marcacao \\ \hline
                    21 & \marcacao & \marcacao & \marcacao & \marcacao \\ \hline
                    22 & \marcacao & \marcacao & \marcacao & \marcacao \\ \hline
                    23 & \marcacao & \marcacao & \marcacao & \marcacao \\ \hline
                    24 & \marcacao & \marcacao & \marcacao & \marcacao \\ \hline
                    25 & \marcacao & \marcacao & \marcacao & \marcacao \\ \hline
                \end{tabular}
            };
            
            % --- Marcadores OMR ---
            \fill[black] ([shift={(-7mm,7mm)}]table.north west) rectangle ++(5mm,-5mm);
            \fill[black] ([shift={(7mm,7mm)}]table.north east) rectangle ++(-5mm,-5mm);
            \fill[black] ([shift={(-7mm,-7mm)}]table.south west) rectangle ++(5mm,5mm);
            \fill[black] ([shift={(7mm,-7mm)}]table.south east) rectangle ++(-5mm,5mm);
        \end{tikzpicture}
    \end{center}
    \renewcommand{\arraystretch}{1.0}
}

% Chamada do gabarito
\Gabarito

\newpage
% --- Bolinha centralizada ---
\newcommand{\marcacao}{%
    \vfill
    \tikz[baseline=(c.base)] \node (c) [circle, draw, minimum width=0.45cm, minimum height=0.45cm] {};
    \vfill
}

% --- Nova coluna centralizada ---
\newcolumntype{C}[1]{>{\centering\arraybackslash}m{#1}}

% --- Cabeçalho ---
\begin{minipage}[t]{0.7\textwidth}
    \textbf{Escola:} CEF04\_PLAN \\
    \textbf{Turma:} 2S - 9º ANO - Matutino \\
    \textbf{Aluno:} HUGO HENRIQUE DA SILVA SOUSA
\end{minipage}
\begin{minipage}[t]{0.29\textwidth}
    \raggedleft
    \qrcode[height=2.2cm]{ 189979-HUGO_HENRIQUE_DA_SILVA_SOUSA-2S }
\end{minipage}

\vspace{0.5cm}
\hrule
\vspace{0.5cm}

% --- Gabarito com 25 questões compactado ---
\newcommand{\Gabarito}{
    \renewcommand{\arraystretch}{1.2} % altura das linhas compacta
    \begin{center}
        \begin{tikzpicture}
            \node (table) [inner sep=0pt, anchor=north west] {
                \begin{tabular}{|C{0.6cm}|C{0.8cm}|C{0.8cm}|C{0.8cm}|C{0.8cm}|}
                    \hline
                    \adjustbox{max width=0.5cm}{\textbf{Nº}} & \textbf{A} & \textbf{B} & \textbf{C} & \textbf{D} \\
                    \hline
                    01 & \marcacao & \marcacao & \marcacao & \marcacao \\ \hline
                    02 & \marcacao & \marcacao & \marcacao & \marcacao \\ \hline
                    03 & \marcacao & \marcacao & \marcacao & \marcacao \\ \hline
                    04 & \marcacao & \marcacao & \marcacao & \marcacao \\ \hline
                    05 & \marcacao & \marcacao & \marcacao & \marcacao \\ \hline
                    06 & \marcacao & \marcacao & \marcacao & \marcacao \\ \hline
                    07 & \marcacao & \marcacao & \marcacao & \marcacao \\ \hline
                    08 & \marcacao & \marcacao & \marcacao & \marcacao \\ \hline
                    09 & \marcacao & \marcacao & \marcacao & \marcacao \\ \hline
                    10 & \marcacao & \marcacao & \marcacao & \marcacao \\ \hline
                    11 & \marcacao & \marcacao & \marcacao & \marcacao \\ \hline
                    12 & \marcacao & \marcacao & \marcacao & \marcacao \\ \hline
                    13 & \marcacao & \marcacao & \marcacao & \marcacao \\ \hline
                    14 & \marcacao & \marcacao & \marcacao & \marcacao \\ \hline
                    15 & \marcacao & \marcacao & \marcacao & \marcacao \\ \hline
                    16 & \marcacao & \marcacao & \marcacao & \marcacao \\ \hline
                    17 & \marcacao & \marcacao & \marcacao & \marcacao \\ \hline
                    18 & \marcacao & \marcacao & \marcacao & \marcacao \\ \hline
                    19 & \marcacao & \marcacao & \marcacao & \marcacao \\ \hline
                    20 & \marcacao & \marcacao & \marcacao & \marcacao \\ \hline
                    21 & \marcacao & \marcacao & \marcacao & \marcacao \\ \hline
                    22 & \marcacao & \marcacao & \marcacao & \marcacao \\ \hline
                    23 & \marcacao & \marcacao & \marcacao & \marcacao \\ \hline
                    24 & \marcacao & \marcacao & \marcacao & \marcacao \\ \hline
                    25 & \marcacao & \marcacao & \marcacao & \marcacao \\ \hline
                \end{tabular}
            };
            
            % --- Marcadores OMR ---
            \fill[black] ([shift={(-7mm,7mm)}]table.north west) rectangle ++(5mm,-5mm);
            \fill[black] ([shift={(7mm,7mm)}]table.north east) rectangle ++(-5mm,-5mm);
            \fill[black] ([shift={(-7mm,-7mm)}]table.south west) rectangle ++(5mm,5mm);
            \fill[black] ([shift={(7mm,-7mm)}]table.south east) rectangle ++(-5mm,5mm);
        \end{tikzpicture}
    \end{center}
    \renewcommand{\arraystretch}{1.0}
}

% Chamada do gabarito
\Gabarito

\newpage
% --- Bolinha centralizada ---
\newcommand{\marcacao}{%
    \vfill
    \tikz[baseline=(c.base)] \node (c) [circle, draw, minimum width=0.45cm, minimum height=0.45cm] {};
    \vfill
}

% --- Nova coluna centralizada ---
\newcolumntype{C}[1]{>{\centering\arraybackslash}m{#1}}

% --- Cabeçalho ---
\begin{minipage}[t]{0.7\textwidth}
    \textbf{Escola:} CEF04\_PLAN \\
    \textbf{Turma:} 2S - 9º ANO - Matutino \\
    \textbf{Aluno:} ISADORA DA CRUZ COSTA
\end{minipage}
\begin{minipage}[t]{0.29\textwidth}
    \raggedleft
    \qrcode[height=2.2cm]{ 536414-ISADORA_DA_CRUZ_COSTA-2S }
\end{minipage}

\vspace{0.5cm}
\hrule
\vspace{0.5cm}

% --- Gabarito com 25 questões compactado ---
\newcommand{\Gabarito}{
    \renewcommand{\arraystretch}{1.2} % altura das linhas compacta
    \begin{center}
        \begin{tikzpicture}
            \node (table) [inner sep=0pt, anchor=north west] {
                \begin{tabular}{|C{0.6cm}|C{0.8cm}|C{0.8cm}|C{0.8cm}|C{0.8cm}|}
                    \hline
                    \adjustbox{max width=0.5cm}{\textbf{Nº}} & \textbf{A} & \textbf{B} & \textbf{C} & \textbf{D} \\
                    \hline
                    01 & \marcacao & \marcacao & \marcacao & \marcacao \\ \hline
                    02 & \marcacao & \marcacao & \marcacao & \marcacao \\ \hline
                    03 & \marcacao & \marcacao & \marcacao & \marcacao \\ \hline
                    04 & \marcacao & \marcacao & \marcacao & \marcacao \\ \hline
                    05 & \marcacao & \marcacao & \marcacao & \marcacao \\ \hline
                    06 & \marcacao & \marcacao & \marcacao & \marcacao \\ \hline
                    07 & \marcacao & \marcacao & \marcacao & \marcacao \\ \hline
                    08 & \marcacao & \marcacao & \marcacao & \marcacao \\ \hline
                    09 & \marcacao & \marcacao & \marcacao & \marcacao \\ \hline
                    10 & \marcacao & \marcacao & \marcacao & \marcacao \\ \hline
                    11 & \marcacao & \marcacao & \marcacao & \marcacao \\ \hline
                    12 & \marcacao & \marcacao & \marcacao & \marcacao \\ \hline
                    13 & \marcacao & \marcacao & \marcacao & \marcacao \\ \hline
                    14 & \marcacao & \marcacao & \marcacao & \marcacao \\ \hline
                    15 & \marcacao & \marcacao & \marcacao & \marcacao \\ \hline
                    16 & \marcacao & \marcacao & \marcacao & \marcacao \\ \hline
                    17 & \marcacao & \marcacao & \marcacao & \marcacao \\ \hline
                    18 & \marcacao & \marcacao & \marcacao & \marcacao \\ \hline
                    19 & \marcacao & \marcacao & \marcacao & \marcacao \\ \hline
                    20 & \marcacao & \marcacao & \marcacao & \marcacao \\ \hline
                    21 & \marcacao & \marcacao & \marcacao & \marcacao \\ \hline
                    22 & \marcacao & \marcacao & \marcacao & \marcacao \\ \hline
                    23 & \marcacao & \marcacao & \marcacao & \marcacao \\ \hline
                    24 & \marcacao & \marcacao & \marcacao & \marcacao \\ \hline
                    25 & \marcacao & \marcacao & \marcacao & \marcacao \\ \hline
                \end{tabular}
            };
            
            % --- Marcadores OMR ---
            \fill[black] ([shift={(-7mm,7mm)}]table.north west) rectangle ++(5mm,-5mm);
            \fill[black] ([shift={(7mm,7mm)}]table.north east) rectangle ++(-5mm,-5mm);
            \fill[black] ([shift={(-7mm,-7mm)}]table.south west) rectangle ++(5mm,5mm);
            \fill[black] ([shift={(7mm,-7mm)}]table.south east) rectangle ++(-5mm,5mm);
        \end{tikzpicture}
    \end{center}
    \renewcommand{\arraystretch}{1.0}
}

% Chamada do gabarito
\Gabarito

\newpage
% --- Bolinha centralizada ---
\newcommand{\marcacao}{%
    \vfill
    \tikz[baseline=(c.base)] \node (c) [circle, draw, minimum width=0.45cm, minimum height=0.45cm] {};
    \vfill
}

% --- Nova coluna centralizada ---
\newcolumntype{C}[1]{>{\centering\arraybackslash}m{#1}}

% --- Cabeçalho ---
\begin{minipage}[t]{0.7\textwidth}
    \textbf{Escola:} CEF04\_PLAN \\
    \textbf{Turma:} 2S - 9º ANO - Matutino \\
    \textbf{Aluno:} JOÃO AURÉLIO XAVIER DOS SANTOS
\end{minipage}
\begin{minipage}[t]{0.29\textwidth}
    \raggedleft
    \qrcode[height=2.2cm]{ 620634-JOÃO_AURÉLIO_XAVIER_DOS_SANTOS-2S }
\end{minipage}

\vspace{0.5cm}
\hrule
\vspace{0.5cm}

% --- Gabarito com 25 questões compactado ---
\newcommand{\Gabarito}{
    \renewcommand{\arraystretch}{1.2} % altura das linhas compacta
    \begin{center}
        \begin{tikzpicture}
            \node (table) [inner sep=0pt, anchor=north west] {
                \begin{tabular}{|C{0.6cm}|C{0.8cm}|C{0.8cm}|C{0.8cm}|C{0.8cm}|}
                    \hline
                    \adjustbox{max width=0.5cm}{\textbf{Nº}} & \textbf{A} & \textbf{B} & \textbf{C} & \textbf{D} \\
                    \hline
                    01 & \marcacao & \marcacao & \marcacao & \marcacao \\ \hline
                    02 & \marcacao & \marcacao & \marcacao & \marcacao \\ \hline
                    03 & \marcacao & \marcacao & \marcacao & \marcacao \\ \hline
                    04 & \marcacao & \marcacao & \marcacao & \marcacao \\ \hline
                    05 & \marcacao & \marcacao & \marcacao & \marcacao \\ \hline
                    06 & \marcacao & \marcacao & \marcacao & \marcacao \\ \hline
                    07 & \marcacao & \marcacao & \marcacao & \marcacao \\ \hline
                    08 & \marcacao & \marcacao & \marcacao & \marcacao \\ \hline
                    09 & \marcacao & \marcacao & \marcacao & \marcacao \\ \hline
                    10 & \marcacao & \marcacao & \marcacao & \marcacao \\ \hline
                    11 & \marcacao & \marcacao & \marcacao & \marcacao \\ \hline
                    12 & \marcacao & \marcacao & \marcacao & \marcacao \\ \hline
                    13 & \marcacao & \marcacao & \marcacao & \marcacao \\ \hline
                    14 & \marcacao & \marcacao & \marcacao & \marcacao \\ \hline
                    15 & \marcacao & \marcacao & \marcacao & \marcacao \\ \hline
                    16 & \marcacao & \marcacao & \marcacao & \marcacao \\ \hline
                    17 & \marcacao & \marcacao & \marcacao & \marcacao \\ \hline
                    18 & \marcacao & \marcacao & \marcacao & \marcacao \\ \hline
                    19 & \marcacao & \marcacao & \marcacao & \marcacao \\ \hline
                    20 & \marcacao & \marcacao & \marcacao & \marcacao \\ \hline
                    21 & \marcacao & \marcacao & \marcacao & \marcacao \\ \hline
                    22 & \marcacao & \marcacao & \marcacao & \marcacao \\ \hline
                    23 & \marcacao & \marcacao & \marcacao & \marcacao \\ \hline
                    24 & \marcacao & \marcacao & \marcacao & \marcacao \\ \hline
                    25 & \marcacao & \marcacao & \marcacao & \marcacao \\ \hline
                \end{tabular}
            };
            
            % --- Marcadores OMR ---
            \fill[black] ([shift={(-7mm,7mm)}]table.north west) rectangle ++(5mm,-5mm);
            \fill[black] ([shift={(7mm,7mm)}]table.north east) rectangle ++(-5mm,-5mm);
            \fill[black] ([shift={(-7mm,-7mm)}]table.south west) rectangle ++(5mm,5mm);
            \fill[black] ([shift={(7mm,-7mm)}]table.south east) rectangle ++(-5mm,5mm);
        \end{tikzpicture}
    \end{center}
    \renewcommand{\arraystretch}{1.0}
}

% Chamada do gabarito
\Gabarito

\newpage
% --- Bolinha centralizada ---
\newcommand{\marcacao}{%
    \vfill
    \tikz[baseline=(c.base)] \node (c) [circle, draw, minimum width=0.45cm, minimum height=0.45cm] {};
    \vfill
}

% --- Nova coluna centralizada ---
\newcolumntype{C}[1]{>{\centering\arraybackslash}m{#1}}

% --- Cabeçalho ---
\begin{minipage}[t]{0.7\textwidth}
    \textbf{Escola:} CEF04\_PLAN \\
    \textbf{Turma:} 2S - 9º ANO - Matutino \\
    \textbf{Aluno:} JOÃO LEONARDO ROCHA GUEDES BONIFÁCIO
\end{minipage}
\begin{minipage}[t]{0.29\textwidth}
    \raggedleft
    \qrcode[height=2.2cm]{ 531404-JOÃO_LEONARDO_ROCHA_GUEDES_BONIFÁCIO-2S }
\end{minipage}

\vspace{0.5cm}
\hrule
\vspace{0.5cm}

% --- Gabarito com 25 questões compactado ---
\newcommand{\Gabarito}{
    \renewcommand{\arraystretch}{1.2} % altura das linhas compacta
    \begin{center}
        \begin{tikzpicture}
            \node (table) [inner sep=0pt, anchor=north west] {
                \begin{tabular}{|C{0.6cm}|C{0.8cm}|C{0.8cm}|C{0.8cm}|C{0.8cm}|}
                    \hline
                    \adjustbox{max width=0.5cm}{\textbf{Nº}} & \textbf{A} & \textbf{B} & \textbf{C} & \textbf{D} \\
                    \hline
                    01 & \marcacao & \marcacao & \marcacao & \marcacao \\ \hline
                    02 & \marcacao & \marcacao & \marcacao & \marcacao \\ \hline
                    03 & \marcacao & \marcacao & \marcacao & \marcacao \\ \hline
                    04 & \marcacao & \marcacao & \marcacao & \marcacao \\ \hline
                    05 & \marcacao & \marcacao & \marcacao & \marcacao \\ \hline
                    06 & \marcacao & \marcacao & \marcacao & \marcacao \\ \hline
                    07 & \marcacao & \marcacao & \marcacao & \marcacao \\ \hline
                    08 & \marcacao & \marcacao & \marcacao & \marcacao \\ \hline
                    09 & \marcacao & \marcacao & \marcacao & \marcacao \\ \hline
                    10 & \marcacao & \marcacao & \marcacao & \marcacao \\ \hline
                    11 & \marcacao & \marcacao & \marcacao & \marcacao \\ \hline
                    12 & \marcacao & \marcacao & \marcacao & \marcacao \\ \hline
                    13 & \marcacao & \marcacao & \marcacao & \marcacao \\ \hline
                    14 & \marcacao & \marcacao & \marcacao & \marcacao \\ \hline
                    15 & \marcacao & \marcacao & \marcacao & \marcacao \\ \hline
                    16 & \marcacao & \marcacao & \marcacao & \marcacao \\ \hline
                    17 & \marcacao & \marcacao & \marcacao & \marcacao \\ \hline
                    18 & \marcacao & \marcacao & \marcacao & \marcacao \\ \hline
                    19 & \marcacao & \marcacao & \marcacao & \marcacao \\ \hline
                    20 & \marcacao & \marcacao & \marcacao & \marcacao \\ \hline
                    21 & \marcacao & \marcacao & \marcacao & \marcacao \\ \hline
                    22 & \marcacao & \marcacao & \marcacao & \marcacao \\ \hline
                    23 & \marcacao & \marcacao & \marcacao & \marcacao \\ \hline
                    24 & \marcacao & \marcacao & \marcacao & \marcacao \\ \hline
                    25 & \marcacao & \marcacao & \marcacao & \marcacao \\ \hline
                \end{tabular}
            };
            
            % --- Marcadores OMR ---
            \fill[black] ([shift={(-7mm,7mm)}]table.north west) rectangle ++(5mm,-5mm);
            \fill[black] ([shift={(7mm,7mm)}]table.north east) rectangle ++(-5mm,-5mm);
            \fill[black] ([shift={(-7mm,-7mm)}]table.south west) rectangle ++(5mm,5mm);
            \fill[black] ([shift={(7mm,-7mm)}]table.south east) rectangle ++(-5mm,5mm);
        \end{tikzpicture}
    \end{center}
    \renewcommand{\arraystretch}{1.0}
}

% Chamada do gabarito
\Gabarito

\newpage
% --- Bolinha centralizada ---
\newcommand{\marcacao}{%
    \vfill
    \tikz[baseline=(c.base)] \node (c) [circle, draw, minimum width=0.45cm, minimum height=0.45cm] {};
    \vfill
}

% --- Nova coluna centralizada ---
\newcolumntype{C}[1]{>{\centering\arraybackslash}m{#1}}

% --- Cabeçalho ---
\begin{minipage}[t]{0.7\textwidth}
    \textbf{Escola:} CEF04\_PLAN \\
    \textbf{Turma:} 2S - 9º ANO - Matutino \\
    \textbf{Aluno:} JOÃO VICTOR LIMA DE PAULA
\end{minipage}
\begin{minipage}[t]{0.29\textwidth}
    \raggedleft
    \qrcode[height=2.2cm]{ 594339-JOÃO_VICTOR_LIMA_DE_PAULA-2S }
\end{minipage}

\vspace{0.5cm}
\hrule
\vspace{0.5cm}

% --- Gabarito com 25 questões compactado ---
\newcommand{\Gabarito}{
    \renewcommand{\arraystretch}{1.2} % altura das linhas compacta
    \begin{center}
        \begin{tikzpicture}
            \node (table) [inner sep=0pt, anchor=north west] {
                \begin{tabular}{|C{0.6cm}|C{0.8cm}|C{0.8cm}|C{0.8cm}|C{0.8cm}|}
                    \hline
                    \adjustbox{max width=0.5cm}{\textbf{Nº}} & \textbf{A} & \textbf{B} & \textbf{C} & \textbf{D} \\
                    \hline
                    01 & \marcacao & \marcacao & \marcacao & \marcacao \\ \hline
                    02 & \marcacao & \marcacao & \marcacao & \marcacao \\ \hline
                    03 & \marcacao & \marcacao & \marcacao & \marcacao \\ \hline
                    04 & \marcacao & \marcacao & \marcacao & \marcacao \\ \hline
                    05 & \marcacao & \marcacao & \marcacao & \marcacao \\ \hline
                    06 & \marcacao & \marcacao & \marcacao & \marcacao \\ \hline
                    07 & \marcacao & \marcacao & \marcacao & \marcacao \\ \hline
                    08 & \marcacao & \marcacao & \marcacao & \marcacao \\ \hline
                    09 & \marcacao & \marcacao & \marcacao & \marcacao \\ \hline
                    10 & \marcacao & \marcacao & \marcacao & \marcacao \\ \hline
                    11 & \marcacao & \marcacao & \marcacao & \marcacao \\ \hline
                    12 & \marcacao & \marcacao & \marcacao & \marcacao \\ \hline
                    13 & \marcacao & \marcacao & \marcacao & \marcacao \\ \hline
                    14 & \marcacao & \marcacao & \marcacao & \marcacao \\ \hline
                    15 & \marcacao & \marcacao & \marcacao & \marcacao \\ \hline
                    16 & \marcacao & \marcacao & \marcacao & \marcacao \\ \hline
                    17 & \marcacao & \marcacao & \marcacao & \marcacao \\ \hline
                    18 & \marcacao & \marcacao & \marcacao & \marcacao \\ \hline
                    19 & \marcacao & \marcacao & \marcacao & \marcacao \\ \hline
                    20 & \marcacao & \marcacao & \marcacao & \marcacao \\ \hline
                    21 & \marcacao & \marcacao & \marcacao & \marcacao \\ \hline
                    22 & \marcacao & \marcacao & \marcacao & \marcacao \\ \hline
                    23 & \marcacao & \marcacao & \marcacao & \marcacao \\ \hline
                    24 & \marcacao & \marcacao & \marcacao & \marcacao \\ \hline
                    25 & \marcacao & \marcacao & \marcacao & \marcacao \\ \hline
                \end{tabular}
            };
            
            % --- Marcadores OMR ---
            \fill[black] ([shift={(-7mm,7mm)}]table.north west) rectangle ++(5mm,-5mm);
            \fill[black] ([shift={(7mm,7mm)}]table.north east) rectangle ++(-5mm,-5mm);
            \fill[black] ([shift={(-7mm,-7mm)}]table.south west) rectangle ++(5mm,5mm);
            \fill[black] ([shift={(7mm,-7mm)}]table.south east) rectangle ++(-5mm,5mm);
        \end{tikzpicture}
    \end{center}
    \renewcommand{\arraystretch}{1.0}
}

% Chamada do gabarito
\Gabarito

\newpage
% --- Bolinha centralizada ---
\newcommand{\marcacao}{%
    \vfill
    \tikz[baseline=(c.base)] \node (c) [circle, draw, minimum width=0.45cm, minimum height=0.45cm] {};
    \vfill
}

% --- Nova coluna centralizada ---
\newcolumntype{C}[1]{>{\centering\arraybackslash}m{#1}}

% --- Cabeçalho ---
\begin{minipage}[t]{0.7\textwidth}
    \textbf{Escola:} CEF04\_PLAN \\
    \textbf{Turma:} 2S - 9º ANO - Matutino \\
    \textbf{Aluno:} JULLYE FIGUEIREDO GOLBERTO
\end{minipage}
\begin{minipage}[t]{0.29\textwidth}
    \raggedleft
    \qrcode[height=2.2cm]{ 20398-JULLYE_FIGUEIREDO_GOLBERTO-2S }
\end{minipage}

\vspace{0.5cm}
\hrule
\vspace{0.5cm}

% --- Gabarito com 25 questões compactado ---
\newcommand{\Gabarito}{
    \renewcommand{\arraystretch}{1.2} % altura das linhas compacta
    \begin{center}
        \begin{tikzpicture}
            \node (table) [inner sep=0pt, anchor=north west] {
                \begin{tabular}{|C{0.6cm}|C{0.8cm}|C{0.8cm}|C{0.8cm}|C{0.8cm}|}
                    \hline
                    \adjustbox{max width=0.5cm}{\textbf{Nº}} & \textbf{A} & \textbf{B} & \textbf{C} & \textbf{D} \\
                    \hline
                    01 & \marcacao & \marcacao & \marcacao & \marcacao \\ \hline
                    02 & \marcacao & \marcacao & \marcacao & \marcacao \\ \hline
                    03 & \marcacao & \marcacao & \marcacao & \marcacao \\ \hline
                    04 & \marcacao & \marcacao & \marcacao & \marcacao \\ \hline
                    05 & \marcacao & \marcacao & \marcacao & \marcacao \\ \hline
                    06 & \marcacao & \marcacao & \marcacao & \marcacao \\ \hline
                    07 & \marcacao & \marcacao & \marcacao & \marcacao \\ \hline
                    08 & \marcacao & \marcacao & \marcacao & \marcacao \\ \hline
                    09 & \marcacao & \marcacao & \marcacao & \marcacao \\ \hline
                    10 & \marcacao & \marcacao & \marcacao & \marcacao \\ \hline
                    11 & \marcacao & \marcacao & \marcacao & \marcacao \\ \hline
                    12 & \marcacao & \marcacao & \marcacao & \marcacao \\ \hline
                    13 & \marcacao & \marcacao & \marcacao & \marcacao \\ \hline
                    14 & \marcacao & \marcacao & \marcacao & \marcacao \\ \hline
                    15 & \marcacao & \marcacao & \marcacao & \marcacao \\ \hline
                    16 & \marcacao & \marcacao & \marcacao & \marcacao \\ \hline
                    17 & \marcacao & \marcacao & \marcacao & \marcacao \\ \hline
                    18 & \marcacao & \marcacao & \marcacao & \marcacao \\ \hline
                    19 & \marcacao & \marcacao & \marcacao & \marcacao \\ \hline
                    20 & \marcacao & \marcacao & \marcacao & \marcacao \\ \hline
                    21 & \marcacao & \marcacao & \marcacao & \marcacao \\ \hline
                    22 & \marcacao & \marcacao & \marcacao & \marcacao \\ \hline
                    23 & \marcacao & \marcacao & \marcacao & \marcacao \\ \hline
                    24 & \marcacao & \marcacao & \marcacao & \marcacao \\ \hline
                    25 & \marcacao & \marcacao & \marcacao & \marcacao \\ \hline
                \end{tabular}
            };
            
            % --- Marcadores OMR ---
            \fill[black] ([shift={(-7mm,7mm)}]table.north west) rectangle ++(5mm,-5mm);
            \fill[black] ([shift={(7mm,7mm)}]table.north east) rectangle ++(-5mm,-5mm);
            \fill[black] ([shift={(-7mm,-7mm)}]table.south west) rectangle ++(5mm,5mm);
            \fill[black] ([shift={(7mm,-7mm)}]table.south east) rectangle ++(-5mm,5mm);
        \end{tikzpicture}
    \end{center}
    \renewcommand{\arraystretch}{1.0}
}

% Chamada do gabarito
\Gabarito

\newpage
% --- Bolinha centralizada ---
\newcommand{\marcacao}{%
    \vfill
    \tikz[baseline=(c.base)] \node (c) [circle, draw, minimum width=0.45cm, minimum height=0.45cm] {};
    \vfill
}

% --- Nova coluna centralizada ---
\newcolumntype{C}[1]{>{\centering\arraybackslash}m{#1}}

% --- Cabeçalho ---
\begin{minipage}[t]{0.7\textwidth}
    \textbf{Escola:} CEF04\_PLAN \\
    \textbf{Turma:} 2S - 9º ANO - Matutino \\
    \textbf{Aluno:} LUÍS HENRIQUE DA SILVA MOTA
\end{minipage}
\begin{minipage}[t]{0.29\textwidth}
    \raggedleft
    \qrcode[height=2.2cm]{ 95946-LUÍS_HENRIQUE_DA_SILVA_MOTA-2S }
\end{minipage}

\vspace{0.5cm}
\hrule
\vspace{0.5cm}

% --- Gabarito com 25 questões compactado ---
\newcommand{\Gabarito}{
    \renewcommand{\arraystretch}{1.2} % altura das linhas compacta
    \begin{center}
        \begin{tikzpicture}
            \node (table) [inner sep=0pt, anchor=north west] {
                \begin{tabular}{|C{0.6cm}|C{0.8cm}|C{0.8cm}|C{0.8cm}|C{0.8cm}|}
                    \hline
                    \adjustbox{max width=0.5cm}{\textbf{Nº}} & \textbf{A} & \textbf{B} & \textbf{C} & \textbf{D} \\
                    \hline
                    01 & \marcacao & \marcacao & \marcacao & \marcacao \\ \hline
                    02 & \marcacao & \marcacao & \marcacao & \marcacao \\ \hline
                    03 & \marcacao & \marcacao & \marcacao & \marcacao \\ \hline
                    04 & \marcacao & \marcacao & \marcacao & \marcacao \\ \hline
                    05 & \marcacao & \marcacao & \marcacao & \marcacao \\ \hline
                    06 & \marcacao & \marcacao & \marcacao & \marcacao \\ \hline
                    07 & \marcacao & \marcacao & \marcacao & \marcacao \\ \hline
                    08 & \marcacao & \marcacao & \marcacao & \marcacao \\ \hline
                    09 & \marcacao & \marcacao & \marcacao & \marcacao \\ \hline
                    10 & \marcacao & \marcacao & \marcacao & \marcacao \\ \hline
                    11 & \marcacao & \marcacao & \marcacao & \marcacao \\ \hline
                    12 & \marcacao & \marcacao & \marcacao & \marcacao \\ \hline
                    13 & \marcacao & \marcacao & \marcacao & \marcacao \\ \hline
                    14 & \marcacao & \marcacao & \marcacao & \marcacao \\ \hline
                    15 & \marcacao & \marcacao & \marcacao & \marcacao \\ \hline
                    16 & \marcacao & \marcacao & \marcacao & \marcacao \\ \hline
                    17 & \marcacao & \marcacao & \marcacao & \marcacao \\ \hline
                    18 & \marcacao & \marcacao & \marcacao & \marcacao \\ \hline
                    19 & \marcacao & \marcacao & \marcacao & \marcacao \\ \hline
                    20 & \marcacao & \marcacao & \marcacao & \marcacao \\ \hline
                    21 & \marcacao & \marcacao & \marcacao & \marcacao \\ \hline
                    22 & \marcacao & \marcacao & \marcacao & \marcacao \\ \hline
                    23 & \marcacao & \marcacao & \marcacao & \marcacao \\ \hline
                    24 & \marcacao & \marcacao & \marcacao & \marcacao \\ \hline
                    25 & \marcacao & \marcacao & \marcacao & \marcacao \\ \hline
                \end{tabular}
            };
            
            % --- Marcadores OMR ---
            \fill[black] ([shift={(-7mm,7mm)}]table.north west) rectangle ++(5mm,-5mm);
            \fill[black] ([shift={(7mm,7mm)}]table.north east) rectangle ++(-5mm,-5mm);
            \fill[black] ([shift={(-7mm,-7mm)}]table.south west) rectangle ++(5mm,5mm);
            \fill[black] ([shift={(7mm,-7mm)}]table.south east) rectangle ++(-5mm,5mm);
        \end{tikzpicture}
    \end{center}
    \renewcommand{\arraystretch}{1.0}
}

% Chamada do gabarito
\Gabarito

\newpage
% --- Bolinha centralizada ---
\newcommand{\marcacao}{%
    \vfill
    \tikz[baseline=(c.base)] \node (c) [circle, draw, minimum width=0.45cm, minimum height=0.45cm] {};
    \vfill
}

% --- Nova coluna centralizada ---
\newcolumntype{C}[1]{>{\centering\arraybackslash}m{#1}}

% --- Cabeçalho ---
\begin{minipage}[t]{0.7\textwidth}
    \textbf{Escola:} CEF04\_PLAN \\
    \textbf{Turma:} 2S - 9º ANO - Matutino \\
    \textbf{Aluno:} MARIA CLARA ALVES FERNANDES
\end{minipage}
\begin{minipage}[t]{0.29\textwidth}
    \raggedleft
    \qrcode[height=2.2cm]{ 592648-MARIA_CLARA_ALVES_FERNANDES-2S }
\end{minipage}

\vspace{0.5cm}
\hrule
\vspace{0.5cm}

% --- Gabarito com 25 questões compactado ---
\newcommand{\Gabarito}{
    \renewcommand{\arraystretch}{1.2} % altura das linhas compacta
    \begin{center}
        \begin{tikzpicture}
            \node (table) [inner sep=0pt, anchor=north west] {
                \begin{tabular}{|C{0.6cm}|C{0.8cm}|C{0.8cm}|C{0.8cm}|C{0.8cm}|}
                    \hline
                    \adjustbox{max width=0.5cm}{\textbf{Nº}} & \textbf{A} & \textbf{B} & \textbf{C} & \textbf{D} \\
                    \hline
                    01 & \marcacao & \marcacao & \marcacao & \marcacao \\ \hline
                    02 & \marcacao & \marcacao & \marcacao & \marcacao \\ \hline
                    03 & \marcacao & \marcacao & \marcacao & \marcacao \\ \hline
                    04 & \marcacao & \marcacao & \marcacao & \marcacao \\ \hline
                    05 & \marcacao & \marcacao & \marcacao & \marcacao \\ \hline
                    06 & \marcacao & \marcacao & \marcacao & \marcacao \\ \hline
                    07 & \marcacao & \marcacao & \marcacao & \marcacao \\ \hline
                    08 & \marcacao & \marcacao & \marcacao & \marcacao \\ \hline
                    09 & \marcacao & \marcacao & \marcacao & \marcacao \\ \hline
                    10 & \marcacao & \marcacao & \marcacao & \marcacao \\ \hline
                    11 & \marcacao & \marcacao & \marcacao & \marcacao \\ \hline
                    12 & \marcacao & \marcacao & \marcacao & \marcacao \\ \hline
                    13 & \marcacao & \marcacao & \marcacao & \marcacao \\ \hline
                    14 & \marcacao & \marcacao & \marcacao & \marcacao \\ \hline
                    15 & \marcacao & \marcacao & \marcacao & \marcacao \\ \hline
                    16 & \marcacao & \marcacao & \marcacao & \marcacao \\ \hline
                    17 & \marcacao & \marcacao & \marcacao & \marcacao \\ \hline
                    18 & \marcacao & \marcacao & \marcacao & \marcacao \\ \hline
                    19 & \marcacao & \marcacao & \marcacao & \marcacao \\ \hline
                    20 & \marcacao & \marcacao & \marcacao & \marcacao \\ \hline
                    21 & \marcacao & \marcacao & \marcacao & \marcacao \\ \hline
                    22 & \marcacao & \marcacao & \marcacao & \marcacao \\ \hline
                    23 & \marcacao & \marcacao & \marcacao & \marcacao \\ \hline
                    24 & \marcacao & \marcacao & \marcacao & \marcacao \\ \hline
                    25 & \marcacao & \marcacao & \marcacao & \marcacao \\ \hline
                \end{tabular}
            };
            
            % --- Marcadores OMR ---
            \fill[black] ([shift={(-7mm,7mm)}]table.north west) rectangle ++(5mm,-5mm);
            \fill[black] ([shift={(7mm,7mm)}]table.north east) rectangle ++(-5mm,-5mm);
            \fill[black] ([shift={(-7mm,-7mm)}]table.south west) rectangle ++(5mm,5mm);
            \fill[black] ([shift={(7mm,-7mm)}]table.south east) rectangle ++(-5mm,5mm);
        \end{tikzpicture}
    \end{center}
    \renewcommand{\arraystretch}{1.0}
}

% Chamada do gabarito
\Gabarito

\newpage
% --- Bolinha centralizada ---
\newcommand{\marcacao}{%
    \vfill
    \tikz[baseline=(c.base)] \node (c) [circle, draw, minimum width=0.45cm, minimum height=0.45cm] {};
    \vfill
}

% --- Nova coluna centralizada ---
\newcolumntype{C}[1]{>{\centering\arraybackslash}m{#1}}

% --- Cabeçalho ---
\begin{minipage}[t]{0.7\textwidth}
    \textbf{Escola:} CEF04\_PLAN \\
    \textbf{Turma:} 2S - 9º ANO - Matutino \\
    \textbf{Aluno:} MARIA GABRIELLY GOMES MARQUES
\end{minipage}
\begin{minipage}[t]{0.29\textwidth}
    \raggedleft
    \qrcode[height=2.2cm]{ 344988-MARIA_GABRIELLY_GOMES_MARQUES-2S }
\end{minipage}

\vspace{0.5cm}
\hrule
\vspace{0.5cm}

% --- Gabarito com 25 questões compactado ---
\newcommand{\Gabarito}{
    \renewcommand{\arraystretch}{1.2} % altura das linhas compacta
    \begin{center}
        \begin{tikzpicture}
            \node (table) [inner sep=0pt, anchor=north west] {
                \begin{tabular}{|C{0.6cm}|C{0.8cm}|C{0.8cm}|C{0.8cm}|C{0.8cm}|}
                    \hline
                    \adjustbox{max width=0.5cm}{\textbf{Nº}} & \textbf{A} & \textbf{B} & \textbf{C} & \textbf{D} \\
                    \hline
                    01 & \marcacao & \marcacao & \marcacao & \marcacao \\ \hline
                    02 & \marcacao & \marcacao & \marcacao & \marcacao \\ \hline
                    03 & \marcacao & \marcacao & \marcacao & \marcacao \\ \hline
                    04 & \marcacao & \marcacao & \marcacao & \marcacao \\ \hline
                    05 & \marcacao & \marcacao & \marcacao & \marcacao \\ \hline
                    06 & \marcacao & \marcacao & \marcacao & \marcacao \\ \hline
                    07 & \marcacao & \marcacao & \marcacao & \marcacao \\ \hline
                    08 & \marcacao & \marcacao & \marcacao & \marcacao \\ \hline
                    09 & \marcacao & \marcacao & \marcacao & \marcacao \\ \hline
                    10 & \marcacao & \marcacao & \marcacao & \marcacao \\ \hline
                    11 & \marcacao & \marcacao & \marcacao & \marcacao \\ \hline
                    12 & \marcacao & \marcacao & \marcacao & \marcacao \\ \hline
                    13 & \marcacao & \marcacao & \marcacao & \marcacao \\ \hline
                    14 & \marcacao & \marcacao & \marcacao & \marcacao \\ \hline
                    15 & \marcacao & \marcacao & \marcacao & \marcacao \\ \hline
                    16 & \marcacao & \marcacao & \marcacao & \marcacao \\ \hline
                    17 & \marcacao & \marcacao & \marcacao & \marcacao \\ \hline
                    18 & \marcacao & \marcacao & \marcacao & \marcacao \\ \hline
                    19 & \marcacao & \marcacao & \marcacao & \marcacao \\ \hline
                    20 & \marcacao & \marcacao & \marcacao & \marcacao \\ \hline
                    21 & \marcacao & \marcacao & \marcacao & \marcacao \\ \hline
                    22 & \marcacao & \marcacao & \marcacao & \marcacao \\ \hline
                    23 & \marcacao & \marcacao & \marcacao & \marcacao \\ \hline
                    24 & \marcacao & \marcacao & \marcacao & \marcacao \\ \hline
                    25 & \marcacao & \marcacao & \marcacao & \marcacao \\ \hline
                \end{tabular}
            };
            
            % --- Marcadores OMR ---
            \fill[black] ([shift={(-7mm,7mm)}]table.north west) rectangle ++(5mm,-5mm);
            \fill[black] ([shift={(7mm,7mm)}]table.north east) rectangle ++(-5mm,-5mm);
            \fill[black] ([shift={(-7mm,-7mm)}]table.south west) rectangle ++(5mm,5mm);
            \fill[black] ([shift={(7mm,-7mm)}]table.south east) rectangle ++(-5mm,5mm);
        \end{tikzpicture}
    \end{center}
    \renewcommand{\arraystretch}{1.0}
}

% Chamada do gabarito
\Gabarito

\newpage
% --- Bolinha centralizada ---
\newcommand{\marcacao}{%
    \vfill
    \tikz[baseline=(c.base)] \node (c) [circle, draw, minimum width=0.45cm, minimum height=0.45cm] {};
    \vfill
}

% --- Nova coluna centralizada ---
\newcolumntype{C}[1]{>{\centering\arraybackslash}m{#1}}

% --- Cabeçalho ---
\begin{minipage}[t]{0.7\textwidth}
    \textbf{Escola:} CEF04\_PLAN \\
    \textbf{Turma:} 2S - 9º ANO - Matutino \\
    \textbf{Aluno:} MARIANA MARTINS DOS SANTOS
\end{minipage}
\begin{minipage}[t]{0.29\textwidth}
    \raggedleft
    \qrcode[height=2.2cm]{ 532827-MARIANA_MARTINS_DOS_SANTOS-2S }
\end{minipage}

\vspace{0.5cm}
\hrule
\vspace{0.5cm}

% --- Gabarito com 25 questões compactado ---
\newcommand{\Gabarito}{
    \renewcommand{\arraystretch}{1.2} % altura das linhas compacta
    \begin{center}
        \begin{tikzpicture}
            \node (table) [inner sep=0pt, anchor=north west] {
                \begin{tabular}{|C{0.6cm}|C{0.8cm}|C{0.8cm}|C{0.8cm}|C{0.8cm}|}
                    \hline
                    \adjustbox{max width=0.5cm}{\textbf{Nº}} & \textbf{A} & \textbf{B} & \textbf{C} & \textbf{D} \\
                    \hline
                    01 & \marcacao & \marcacao & \marcacao & \marcacao \\ \hline
                    02 & \marcacao & \marcacao & \marcacao & \marcacao \\ \hline
                    03 & \marcacao & \marcacao & \marcacao & \marcacao \\ \hline
                    04 & \marcacao & \marcacao & \marcacao & \marcacao \\ \hline
                    05 & \marcacao & \marcacao & \marcacao & \marcacao \\ \hline
                    06 & \marcacao & \marcacao & \marcacao & \marcacao \\ \hline
                    07 & \marcacao & \marcacao & \marcacao & \marcacao \\ \hline
                    08 & \marcacao & \marcacao & \marcacao & \marcacao \\ \hline
                    09 & \marcacao & \marcacao & \marcacao & \marcacao \\ \hline
                    10 & \marcacao & \marcacao & \marcacao & \marcacao \\ \hline
                    11 & \marcacao & \marcacao & \marcacao & \marcacao \\ \hline
                    12 & \marcacao & \marcacao & \marcacao & \marcacao \\ \hline
                    13 & \marcacao & \marcacao & \marcacao & \marcacao \\ \hline
                    14 & \marcacao & \marcacao & \marcacao & \marcacao \\ \hline
                    15 & \marcacao & \marcacao & \marcacao & \marcacao \\ \hline
                    16 & \marcacao & \marcacao & \marcacao & \marcacao \\ \hline
                    17 & \marcacao & \marcacao & \marcacao & \marcacao \\ \hline
                    18 & \marcacao & \marcacao & \marcacao & \marcacao \\ \hline
                    19 & \marcacao & \marcacao & \marcacao & \marcacao \\ \hline
                    20 & \marcacao & \marcacao & \marcacao & \marcacao \\ \hline
                    21 & \marcacao & \marcacao & \marcacao & \marcacao \\ \hline
                    22 & \marcacao & \marcacao & \marcacao & \marcacao \\ \hline
                    23 & \marcacao & \marcacao & \marcacao & \marcacao \\ \hline
                    24 & \marcacao & \marcacao & \marcacao & \marcacao \\ \hline
                    25 & \marcacao & \marcacao & \marcacao & \marcacao \\ \hline
                \end{tabular}
            };
            
            % --- Marcadores OMR ---
            \fill[black] ([shift={(-7mm,7mm)}]table.north west) rectangle ++(5mm,-5mm);
            \fill[black] ([shift={(7mm,7mm)}]table.north east) rectangle ++(-5mm,-5mm);
            \fill[black] ([shift={(-7mm,-7mm)}]table.south west) rectangle ++(5mm,5mm);
            \fill[black] ([shift={(7mm,-7mm)}]table.south east) rectangle ++(-5mm,5mm);
        \end{tikzpicture}
    \end{center}
    \renewcommand{\arraystretch}{1.0}
}

% Chamada do gabarito
\Gabarito

\newpage
% --- Bolinha centralizada ---
\newcommand{\marcacao}{%
    \vfill
    \tikz[baseline=(c.base)] \node (c) [circle, draw, minimum width=0.45cm, minimum height=0.45cm] {};
    \vfill
}

% --- Nova coluna centralizada ---
\newcolumntype{C}[1]{>{\centering\arraybackslash}m{#1}}

% --- Cabeçalho ---
\begin{minipage}[t]{0.7\textwidth}
    \textbf{Escola:} CEF04\_PLAN \\
    \textbf{Turma:} 2S - 9º ANO - Matutino \\
    \textbf{Aluno:} NICOLI OLIVEIRA DE JESUS GUEDES
\end{minipage}
\begin{minipage}[t]{0.29\textwidth}
    \raggedleft
    \qrcode[height=2.2cm]{ 12629-NICOLI_OLIVEIRA_DE_JESUS_GUEDES-2S }
\end{minipage}

\vspace{0.5cm}
\hrule
\vspace{0.5cm}

% --- Gabarito com 25 questões compactado ---
\newcommand{\Gabarito}{
    \renewcommand{\arraystretch}{1.2} % altura das linhas compacta
    \begin{center}
        \begin{tikzpicture}
            \node (table) [inner sep=0pt, anchor=north west] {
                \begin{tabular}{|C{0.6cm}|C{0.8cm}|C{0.8cm}|C{0.8cm}|C{0.8cm}|}
                    \hline
                    \adjustbox{max width=0.5cm}{\textbf{Nº}} & \textbf{A} & \textbf{B} & \textbf{C} & \textbf{D} \\
                    \hline
                    01 & \marcacao & \marcacao & \marcacao & \marcacao \\ \hline
                    02 & \marcacao & \marcacao & \marcacao & \marcacao \\ \hline
                    03 & \marcacao & \marcacao & \marcacao & \marcacao \\ \hline
                    04 & \marcacao & \marcacao & \marcacao & \marcacao \\ \hline
                    05 & \marcacao & \marcacao & \marcacao & \marcacao \\ \hline
                    06 & \marcacao & \marcacao & \marcacao & \marcacao \\ \hline
                    07 & \marcacao & \marcacao & \marcacao & \marcacao \\ \hline
                    08 & \marcacao & \marcacao & \marcacao & \marcacao \\ \hline
                    09 & \marcacao & \marcacao & \marcacao & \marcacao \\ \hline
                    10 & \marcacao & \marcacao & \marcacao & \marcacao \\ \hline
                    11 & \marcacao & \marcacao & \marcacao & \marcacao \\ \hline
                    12 & \marcacao & \marcacao & \marcacao & \marcacao \\ \hline
                    13 & \marcacao & \marcacao & \marcacao & \marcacao \\ \hline
                    14 & \marcacao & \marcacao & \marcacao & \marcacao \\ \hline
                    15 & \marcacao & \marcacao & \marcacao & \marcacao \\ \hline
                    16 & \marcacao & \marcacao & \marcacao & \marcacao \\ \hline
                    17 & \marcacao & \marcacao & \marcacao & \marcacao \\ \hline
                    18 & \marcacao & \marcacao & \marcacao & \marcacao \\ \hline
                    19 & \marcacao & \marcacao & \marcacao & \marcacao \\ \hline
                    20 & \marcacao & \marcacao & \marcacao & \marcacao \\ \hline
                    21 & \marcacao & \marcacao & \marcacao & \marcacao \\ \hline
                    22 & \marcacao & \marcacao & \marcacao & \marcacao \\ \hline
                    23 & \marcacao & \marcacao & \marcacao & \marcacao \\ \hline
                    24 & \marcacao & \marcacao & \marcacao & \marcacao \\ \hline
                    25 & \marcacao & \marcacao & \marcacao & \marcacao \\ \hline
                \end{tabular}
            };
            
            % --- Marcadores OMR ---
            \fill[black] ([shift={(-7mm,7mm)}]table.north west) rectangle ++(5mm,-5mm);
            \fill[black] ([shift={(7mm,7mm)}]table.north east) rectangle ++(-5mm,-5mm);
            \fill[black] ([shift={(-7mm,-7mm)}]table.south west) rectangle ++(5mm,5mm);
            \fill[black] ([shift={(7mm,-7mm)}]table.south east) rectangle ++(-5mm,5mm);
        \end{tikzpicture}
    \end{center}
    \renewcommand{\arraystretch}{1.0}
}

% Chamada do gabarito
\Gabarito

\newpage
% --- Bolinha centralizada ---
\newcommand{\marcacao}{%
    \vfill
    \tikz[baseline=(c.base)] \node (c) [circle, draw, minimum width=0.45cm, minimum height=0.45cm] {};
    \vfill
}

% --- Nova coluna centralizada ---
\newcolumntype{C}[1]{>{\centering\arraybackslash}m{#1}}

% --- Cabeçalho ---
\begin{minipage}[t]{0.7\textwidth}
    \textbf{Escola:} CEF04\_PLAN \\
    \textbf{Turma:} 2S - 9º ANO - Matutino \\
    \textbf{Aluno:} PEDRO HENRIQUE FERREIRA
\end{minipage}
\begin{minipage}[t]{0.29\textwidth}
    \raggedleft
    \qrcode[height=2.2cm]{ 1166304-PEDRO_HENRIQUE_FERREIRA-2S }
\end{minipage}

\vspace{0.5cm}
\hrule
\vspace{0.5cm}

% --- Gabarito com 25 questões compactado ---
\newcommand{\Gabarito}{
    \renewcommand{\arraystretch}{1.2} % altura das linhas compacta
    \begin{center}
        \begin{tikzpicture}
            \node (table) [inner sep=0pt, anchor=north west] {
                \begin{tabular}{|C{0.6cm}|C{0.8cm}|C{0.8cm}|C{0.8cm}|C{0.8cm}|}
                    \hline
                    \adjustbox{max width=0.5cm}{\textbf{Nº}} & \textbf{A} & \textbf{B} & \textbf{C} & \textbf{D} \\
                    \hline
                    01 & \marcacao & \marcacao & \marcacao & \marcacao \\ \hline
                    02 & \marcacao & \marcacao & \marcacao & \marcacao \\ \hline
                    03 & \marcacao & \marcacao & \marcacao & \marcacao \\ \hline
                    04 & \marcacao & \marcacao & \marcacao & \marcacao \\ \hline
                    05 & \marcacao & \marcacao & \marcacao & \marcacao \\ \hline
                    06 & \marcacao & \marcacao & \marcacao & \marcacao \\ \hline
                    07 & \marcacao & \marcacao & \marcacao & \marcacao \\ \hline
                    08 & \marcacao & \marcacao & \marcacao & \marcacao \\ \hline
                    09 & \marcacao & \marcacao & \marcacao & \marcacao \\ \hline
                    10 & \marcacao & \marcacao & \marcacao & \marcacao \\ \hline
                    11 & \marcacao & \marcacao & \marcacao & \marcacao \\ \hline
                    12 & \marcacao & \marcacao & \marcacao & \marcacao \\ \hline
                    13 & \marcacao & \marcacao & \marcacao & \marcacao \\ \hline
                    14 & \marcacao & \marcacao & \marcacao & \marcacao \\ \hline
                    15 & \marcacao & \marcacao & \marcacao & \marcacao \\ \hline
                    16 & \marcacao & \marcacao & \marcacao & \marcacao \\ \hline
                    17 & \marcacao & \marcacao & \marcacao & \marcacao \\ \hline
                    18 & \marcacao & \marcacao & \marcacao & \marcacao \\ \hline
                    19 & \marcacao & \marcacao & \marcacao & \marcacao \\ \hline
                    20 & \marcacao & \marcacao & \marcacao & \marcacao \\ \hline
                    21 & \marcacao & \marcacao & \marcacao & \marcacao \\ \hline
                    22 & \marcacao & \marcacao & \marcacao & \marcacao \\ \hline
                    23 & \marcacao & \marcacao & \marcacao & \marcacao \\ \hline
                    24 & \marcacao & \marcacao & \marcacao & \marcacao \\ \hline
                    25 & \marcacao & \marcacao & \marcacao & \marcacao \\ \hline
                \end{tabular}
            };
            
            % --- Marcadores OMR ---
            \fill[black] ([shift={(-7mm,7mm)}]table.north west) rectangle ++(5mm,-5mm);
            \fill[black] ([shift={(7mm,7mm)}]table.north east) rectangle ++(-5mm,-5mm);
            \fill[black] ([shift={(-7mm,-7mm)}]table.south west) rectangle ++(5mm,5mm);
            \fill[black] ([shift={(7mm,-7mm)}]table.south east) rectangle ++(-5mm,5mm);
        \end{tikzpicture}
    \end{center}
    \renewcommand{\arraystretch}{1.0}
}

% Chamada do gabarito
\Gabarito

\newpage
% --- Bolinha centralizada ---
\newcommand{\marcacao}{%
    \vfill
    \tikz[baseline=(c.base)] \node (c) [circle, draw, minimum width=0.45cm, minimum height=0.45cm] {};
    \vfill
}

% --- Nova coluna centralizada ---
\newcolumntype{C}[1]{>{\centering\arraybackslash}m{#1}}

% --- Cabeçalho ---
\begin{minipage}[t]{0.7\textwidth}
    \textbf{Escola:} CEF04\_PLAN \\
    \textbf{Turma:} 2S - 9º ANO - Matutino \\
    \textbf{Aluno:} ROGER BATISTA RODRIGUES
\end{minipage}
\begin{minipage}[t]{0.29\textwidth}
    \raggedleft
    \qrcode[height=2.2cm]{ 520767-ROGER_BATISTA_RODRIGUES-2S }
\end{minipage}

\vspace{0.5cm}
\hrule
\vspace{0.5cm}

% --- Gabarito com 25 questões compactado ---
\newcommand{\Gabarito}{
    \renewcommand{\arraystretch}{1.2} % altura das linhas compacta
    \begin{center}
        \begin{tikzpicture}
            \node (table) [inner sep=0pt, anchor=north west] {
                \begin{tabular}{|C{0.6cm}|C{0.8cm}|C{0.8cm}|C{0.8cm}|C{0.8cm}|}
                    \hline
                    \adjustbox{max width=0.5cm}{\textbf{Nº}} & \textbf{A} & \textbf{B} & \textbf{C} & \textbf{D} \\
                    \hline
                    01 & \marcacao & \marcacao & \marcacao & \marcacao \\ \hline
                    02 & \marcacao & \marcacao & \marcacao & \marcacao \\ \hline
                    03 & \marcacao & \marcacao & \marcacao & \marcacao \\ \hline
                    04 & \marcacao & \marcacao & \marcacao & \marcacao \\ \hline
                    05 & \marcacao & \marcacao & \marcacao & \marcacao \\ \hline
                    06 & \marcacao & \marcacao & \marcacao & \marcacao \\ \hline
                    07 & \marcacao & \marcacao & \marcacao & \marcacao \\ \hline
                    08 & \marcacao & \marcacao & \marcacao & \marcacao \\ \hline
                    09 & \marcacao & \marcacao & \marcacao & \marcacao \\ \hline
                    10 & \marcacao & \marcacao & \marcacao & \marcacao \\ \hline
                    11 & \marcacao & \marcacao & \marcacao & \marcacao \\ \hline
                    12 & \marcacao & \marcacao & \marcacao & \marcacao \\ \hline
                    13 & \marcacao & \marcacao & \marcacao & \marcacao \\ \hline
                    14 & \marcacao & \marcacao & \marcacao & \marcacao \\ \hline
                    15 & \marcacao & \marcacao & \marcacao & \marcacao \\ \hline
                    16 & \marcacao & \marcacao & \marcacao & \marcacao \\ \hline
                    17 & \marcacao & \marcacao & \marcacao & \marcacao \\ \hline
                    18 & \marcacao & \marcacao & \marcacao & \marcacao \\ \hline
                    19 & \marcacao & \marcacao & \marcacao & \marcacao \\ \hline
                    20 & \marcacao & \marcacao & \marcacao & \marcacao \\ \hline
                    21 & \marcacao & \marcacao & \marcacao & \marcacao \\ \hline
                    22 & \marcacao & \marcacao & \marcacao & \marcacao \\ \hline
                    23 & \marcacao & \marcacao & \marcacao & \marcacao \\ \hline
                    24 & \marcacao & \marcacao & \marcacao & \marcacao \\ \hline
                    25 & \marcacao & \marcacao & \marcacao & \marcacao \\ \hline
                \end{tabular}
            };
            
            % --- Marcadores OMR ---
            \fill[black] ([shift={(-7mm,7mm)}]table.north west) rectangle ++(5mm,-5mm);
            \fill[black] ([shift={(7mm,7mm)}]table.north east) rectangle ++(-5mm,-5mm);
            \fill[black] ([shift={(-7mm,-7mm)}]table.south west) rectangle ++(5mm,5mm);
            \fill[black] ([shift={(7mm,-7mm)}]table.south east) rectangle ++(-5mm,5mm);
        \end{tikzpicture}
    \end{center}
    \renewcommand{\arraystretch}{1.0}
}

% Chamada do gabarito
\Gabarito

\newpage
% --- Bolinha centralizada ---
\newcommand{\marcacao}{%
    \vfill
    \tikz[baseline=(c.base)] \node (c) [circle, draw, minimum width=0.45cm, minimum height=0.45cm] {};
    \vfill
}

% --- Nova coluna centralizada ---
\newcolumntype{C}[1]{>{\centering\arraybackslash}m{#1}}

% --- Cabeçalho ---
\begin{minipage}[t]{0.7\textwidth}
    \textbf{Escola:} CEF04\_PLAN \\
    \textbf{Turma:} 2S - 9º ANO - Matutino \\
    \textbf{Aluno:} ULISSES EMANUEL DA SILVA SOUZA
\end{minipage}
\begin{minipage}[t]{0.29\textwidth}
    \raggedleft
    \qrcode[height=2.2cm]{ 615114-ULISSES_EMANUEL_DA_SILVA_SOUZA-2S }
\end{minipage}

\vspace{0.5cm}
\hrule
\vspace{0.5cm}

% --- Gabarito com 25 questões compactado ---
\newcommand{\Gabarito}{
    \renewcommand{\arraystretch}{1.2} % altura das linhas compacta
    \begin{center}
        \begin{tikzpicture}
            \node (table) [inner sep=0pt, anchor=north west] {
                \begin{tabular}{|C{0.6cm}|C{0.8cm}|C{0.8cm}|C{0.8cm}|C{0.8cm}|}
                    \hline
                    \adjustbox{max width=0.5cm}{\textbf{Nº}} & \textbf{A} & \textbf{B} & \textbf{C} & \textbf{D} \\
                    \hline
                    01 & \marcacao & \marcacao & \marcacao & \marcacao \\ \hline
                    02 & \marcacao & \marcacao & \marcacao & \marcacao \\ \hline
                    03 & \marcacao & \marcacao & \marcacao & \marcacao \\ \hline
                    04 & \marcacao & \marcacao & \marcacao & \marcacao \\ \hline
                    05 & \marcacao & \marcacao & \marcacao & \marcacao \\ \hline
                    06 & \marcacao & \marcacao & \marcacao & \marcacao \\ \hline
                    07 & \marcacao & \marcacao & \marcacao & \marcacao \\ \hline
                    08 & \marcacao & \marcacao & \marcacao & \marcacao \\ \hline
                    09 & \marcacao & \marcacao & \marcacao & \marcacao \\ \hline
                    10 & \marcacao & \marcacao & \marcacao & \marcacao \\ \hline
                    11 & \marcacao & \marcacao & \marcacao & \marcacao \\ \hline
                    12 & \marcacao & \marcacao & \marcacao & \marcacao \\ \hline
                    13 & \marcacao & \marcacao & \marcacao & \marcacao \\ \hline
                    14 & \marcacao & \marcacao & \marcacao & \marcacao \\ \hline
                    15 & \marcacao & \marcacao & \marcacao & \marcacao \\ \hline
                    16 & \marcacao & \marcacao & \marcacao & \marcacao \\ \hline
                    17 & \marcacao & \marcacao & \marcacao & \marcacao \\ \hline
                    18 & \marcacao & \marcacao & \marcacao & \marcacao \\ \hline
                    19 & \marcacao & \marcacao & \marcacao & \marcacao \\ \hline
                    20 & \marcacao & \marcacao & \marcacao & \marcacao \\ \hline
                    21 & \marcacao & \marcacao & \marcacao & \marcacao \\ \hline
                    22 & \marcacao & \marcacao & \marcacao & \marcacao \\ \hline
                    23 & \marcacao & \marcacao & \marcacao & \marcacao \\ \hline
                    24 & \marcacao & \marcacao & \marcacao & \marcacao \\ \hline
                    25 & \marcacao & \marcacao & \marcacao & \marcacao \\ \hline
                \end{tabular}
            };
            
            % --- Marcadores OMR ---
            \fill[black] ([shift={(-7mm,7mm)}]table.north west) rectangle ++(5mm,-5mm);
            \fill[black] ([shift={(7mm,7mm)}]table.north east) rectangle ++(-5mm,-5mm);
            \fill[black] ([shift={(-7mm,-7mm)}]table.south west) rectangle ++(5mm,5mm);
            \fill[black] ([shift={(7mm,-7mm)}]table.south east) rectangle ++(-5mm,5mm);
        \end{tikzpicture}
    \end{center}
    \renewcommand{\arraystretch}{1.0}
}

% Chamada do gabarito
\Gabarito

\newpage
% --- Bolinha centralizada ---
\newcommand{\marcacao}{%
    \vfill
    \tikz[baseline=(c.base)] \node (c) [circle, draw, minimum width=0.45cm, minimum height=0.45cm] {};
    \vfill
}

% --- Nova coluna centralizada ---
\newcolumntype{C}[1]{>{\centering\arraybackslash}m{#1}}

% --- Cabeçalho ---
\begin{minipage}[t]{0.7\textwidth}
    \textbf{Escola:} CEF04\_PLAN \\
    \textbf{Turma:} 2S - 9º ANO - Matutino \\
    \textbf{Aluno:} VINÍCIUS YURI SOARES LUCENA
\end{minipage}
\begin{minipage}[t]{0.29\textwidth}
    \raggedleft
    \qrcode[height=2.2cm]{ 539317-VINÍCIUS_YURI_SOARES_LUCENA-2S }
\end{minipage}

\vspace{0.5cm}
\hrule
\vspace{0.5cm}

% --- Gabarito com 25 questões compactado ---
\newcommand{\Gabarito}{
    \renewcommand{\arraystretch}{1.2} % altura das linhas compacta
    \begin{center}
        \begin{tikzpicture}
            \node (table) [inner sep=0pt, anchor=north west] {
                \begin{tabular}{|C{0.6cm}|C{0.8cm}|C{0.8cm}|C{0.8cm}|C{0.8cm}|}
                    \hline
                    \adjustbox{max width=0.5cm}{\textbf{Nº}} & \textbf{A} & \textbf{B} & \textbf{C} & \textbf{D} \\
                    \hline
                    01 & \marcacao & \marcacao & \marcacao & \marcacao \\ \hline
                    02 & \marcacao & \marcacao & \marcacao & \marcacao \\ \hline
                    03 & \marcacao & \marcacao & \marcacao & \marcacao \\ \hline
                    04 & \marcacao & \marcacao & \marcacao & \marcacao \\ \hline
                    05 & \marcacao & \marcacao & \marcacao & \marcacao \\ \hline
                    06 & \marcacao & \marcacao & \marcacao & \marcacao \\ \hline
                    07 & \marcacao & \marcacao & \marcacao & \marcacao \\ \hline
                    08 & \marcacao & \marcacao & \marcacao & \marcacao \\ \hline
                    09 & \marcacao & \marcacao & \marcacao & \marcacao \\ \hline
                    10 & \marcacao & \marcacao & \marcacao & \marcacao \\ \hline
                    11 & \marcacao & \marcacao & \marcacao & \marcacao \\ \hline
                    12 & \marcacao & \marcacao & \marcacao & \marcacao \\ \hline
                    13 & \marcacao & \marcacao & \marcacao & \marcacao \\ \hline
                    14 & \marcacao & \marcacao & \marcacao & \marcacao \\ \hline
                    15 & \marcacao & \marcacao & \marcacao & \marcacao \\ \hline
                    16 & \marcacao & \marcacao & \marcacao & \marcacao \\ \hline
                    17 & \marcacao & \marcacao & \marcacao & \marcacao \\ \hline
                    18 & \marcacao & \marcacao & \marcacao & \marcacao \\ \hline
                    19 & \marcacao & \marcacao & \marcacao & \marcacao \\ \hline
                    20 & \marcacao & \marcacao & \marcacao & \marcacao \\ \hline
                    21 & \marcacao & \marcacao & \marcacao & \marcacao \\ \hline
                    22 & \marcacao & \marcacao & \marcacao & \marcacao \\ \hline
                    23 & \marcacao & \marcacao & \marcacao & \marcacao \\ \hline
                    24 & \marcacao & \marcacao & \marcacao & \marcacao \\ \hline
                    25 & \marcacao & \marcacao & \marcacao & \marcacao \\ \hline
                \end{tabular}
            };
            
            % --- Marcadores OMR ---
            \fill[black] ([shift={(-7mm,7mm)}]table.north west) rectangle ++(5mm,-5mm);
            \fill[black] ([shift={(7mm,7mm)}]table.north east) rectangle ++(-5mm,-5mm);
            \fill[black] ([shift={(-7mm,-7mm)}]table.south west) rectangle ++(5mm,5mm);
            \fill[black] ([shift={(7mm,-7mm)}]table.south east) rectangle ++(-5mm,5mm);
        \end{tikzpicture}
    \end{center}
    \renewcommand{\arraystretch}{1.0}
}

% Chamada do gabarito
\Gabarito

\newpage
% --- Bolinha centralizada ---
\newcommand{\marcacao}{%
    \vfill
    \tikz[baseline=(c.base)] \node (c) [circle, draw, minimum width=0.45cm, minimum height=0.45cm] {};
    \vfill
}

% --- Nova coluna centralizada ---
\newcolumntype{C}[1]{>{\centering\arraybackslash}m{#1}}

% --- Cabeçalho ---
\begin{minipage}[t]{0.7\textwidth}
    \textbf{Escola:} CEF04\_PLAN \\
    \textbf{Turma:} 2S - 9º ANO - Matutino \\
    \textbf{Aluno:} WAGNER DA SILVA SANTOS
\end{minipage}
\begin{minipage}[t]{0.29\textwidth}
    \raggedleft
    \qrcode[height=2.2cm]{ 496651-WAGNER_DA_SILVA_SANTOS-2S }
\end{minipage}

\vspace{0.5cm}
\hrule
\vspace{0.5cm}

% --- Gabarito com 25 questões compactado ---
\newcommand{\Gabarito}{
    \renewcommand{\arraystretch}{1.2} % altura das linhas compacta
    \begin{center}
        \begin{tikzpicture}
            \node (table) [inner sep=0pt, anchor=north west] {
                \begin{tabular}{|C{0.6cm}|C{0.8cm}|C{0.8cm}|C{0.8cm}|C{0.8cm}|}
                    \hline
                    \adjustbox{max width=0.5cm}{\textbf{Nº}} & \textbf{A} & \textbf{B} & \textbf{C} & \textbf{D} \\
                    \hline
                    01 & \marcacao & \marcacao & \marcacao & \marcacao \\ \hline
                    02 & \marcacao & \marcacao & \marcacao & \marcacao \\ \hline
                    03 & \marcacao & \marcacao & \marcacao & \marcacao \\ \hline
                    04 & \marcacao & \marcacao & \marcacao & \marcacao \\ \hline
                    05 & \marcacao & \marcacao & \marcacao & \marcacao \\ \hline
                    06 & \marcacao & \marcacao & \marcacao & \marcacao \\ \hline
                    07 & \marcacao & \marcacao & \marcacao & \marcacao \\ \hline
                    08 & \marcacao & \marcacao & \marcacao & \marcacao \\ \hline
                    09 & \marcacao & \marcacao & \marcacao & \marcacao \\ \hline
                    10 & \marcacao & \marcacao & \marcacao & \marcacao \\ \hline
                    11 & \marcacao & \marcacao & \marcacao & \marcacao \\ \hline
                    12 & \marcacao & \marcacao & \marcacao & \marcacao \\ \hline
                    13 & \marcacao & \marcacao & \marcacao & \marcacao \\ \hline
                    14 & \marcacao & \marcacao & \marcacao & \marcacao \\ \hline
                    15 & \marcacao & \marcacao & \marcacao & \marcacao \\ \hline
                    16 & \marcacao & \marcacao & \marcacao & \marcacao \\ \hline
                    17 & \marcacao & \marcacao & \marcacao & \marcacao \\ \hline
                    18 & \marcacao & \marcacao & \marcacao & \marcacao \\ \hline
                    19 & \marcacao & \marcacao & \marcacao & \marcacao \\ \hline
                    20 & \marcacao & \marcacao & \marcacao & \marcacao \\ \hline
                    21 & \marcacao & \marcacao & \marcacao & \marcacao \\ \hline
                    22 & \marcacao & \marcacao & \marcacao & \marcacao \\ \hline
                    23 & \marcacao & \marcacao & \marcacao & \marcacao \\ \hline
                    24 & \marcacao & \marcacao & \marcacao & \marcacao \\ \hline
                    25 & \marcacao & \marcacao & \marcacao & \marcacao \\ \hline
                \end{tabular}
            };
            
            % --- Marcadores OMR ---
            \fill[black] ([shift={(-7mm,7mm)}]table.north west) rectangle ++(5mm,-5mm);
            \fill[black] ([shift={(7mm,7mm)}]table.north east) rectangle ++(-5mm,-5mm);
            \fill[black] ([shift={(-7mm,-7mm)}]table.south west) rectangle ++(5mm,5mm);
            \fill[black] ([shift={(7mm,-7mm)}]table.south east) rectangle ++(-5mm,5mm);
        \end{tikzpicture}
    \end{center}
    \renewcommand{\arraystretch}{1.0}
}

% Chamada do gabarito
\Gabarito

\newpage
% --- Bolinha centralizada ---
\newcommand{\marcacao}{%
    \vfill
    \tikz[baseline=(c.base)] \node (c) [circle, draw, minimum width=0.45cm, minimum height=0.45cm] {};
    \vfill
}

% --- Nova coluna centralizada ---
\newcolumntype{C}[1]{>{\centering\arraybackslash}m{#1}}

% --- Cabeçalho ---
\begin{minipage}[t]{0.7\textwidth}
    \textbf{Escola:} CEF04\_PLAN \\
    \textbf{Turma:} 2S - 9º ANO - Matutino \\
    \textbf{Aluno:} YASMIN GABRIELE LILI FÉLIX
\end{minipage}
\begin{minipage}[t]{0.29\textwidth}
    \raggedleft
    \qrcode[height=2.2cm]{ 692331-YASMIN_GABRIELE_LILI_FÉLIX-2S }
\end{minipage}

\vspace{0.5cm}
\hrule
\vspace{0.5cm}

% --- Gabarito com 25 questões compactado ---
\newcommand{\Gabarito}{
    \renewcommand{\arraystretch}{1.2} % altura das linhas compacta
    \begin{center}
        \begin{tikzpicture}
            \node (table) [inner sep=0pt, anchor=north west] {
                \begin{tabular}{|C{0.6cm}|C{0.8cm}|C{0.8cm}|C{0.8cm}|C{0.8cm}|}
                    \hline
                    \adjustbox{max width=0.5cm}{\textbf{Nº}} & \textbf{A} & \textbf{B} & \textbf{C} & \textbf{D} \\
                    \hline
                    01 & \marcacao & \marcacao & \marcacao & \marcacao \\ \hline
                    02 & \marcacao & \marcacao & \marcacao & \marcacao \\ \hline
                    03 & \marcacao & \marcacao & \marcacao & \marcacao \\ \hline
                    04 & \marcacao & \marcacao & \marcacao & \marcacao \\ \hline
                    05 & \marcacao & \marcacao & \marcacao & \marcacao \\ \hline
                    06 & \marcacao & \marcacao & \marcacao & \marcacao \\ \hline
                    07 & \marcacao & \marcacao & \marcacao & \marcacao \\ \hline
                    08 & \marcacao & \marcacao & \marcacao & \marcacao \\ \hline
                    09 & \marcacao & \marcacao & \marcacao & \marcacao \\ \hline
                    10 & \marcacao & \marcacao & \marcacao & \marcacao \\ \hline
                    11 & \marcacao & \marcacao & \marcacao & \marcacao \\ \hline
                    12 & \marcacao & \marcacao & \marcacao & \marcacao \\ \hline
                    13 & \marcacao & \marcacao & \marcacao & \marcacao \\ \hline
                    14 & \marcacao & \marcacao & \marcacao & \marcacao \\ \hline
                    15 & \marcacao & \marcacao & \marcacao & \marcacao \\ \hline
                    16 & \marcacao & \marcacao & \marcacao & \marcacao \\ \hline
                    17 & \marcacao & \marcacao & \marcacao & \marcacao \\ \hline
                    18 & \marcacao & \marcacao & \marcacao & \marcacao \\ \hline
                    19 & \marcacao & \marcacao & \marcacao & \marcacao \\ \hline
                    20 & \marcacao & \marcacao & \marcacao & \marcacao \\ \hline
                    21 & \marcacao & \marcacao & \marcacao & \marcacao \\ \hline
                    22 & \marcacao & \marcacao & \marcacao & \marcacao \\ \hline
                    23 & \marcacao & \marcacao & \marcacao & \marcacao \\ \hline
                    24 & \marcacao & \marcacao & \marcacao & \marcacao \\ \hline
                    25 & \marcacao & \marcacao & \marcacao & \marcacao \\ \hline
                \end{tabular}
            };
            
            % --- Marcadores OMR ---
            \fill[black] ([shift={(-7mm,7mm)}]table.north west) rectangle ++(5mm,-5mm);
            \fill[black] ([shift={(7mm,7mm)}]table.north east) rectangle ++(-5mm,-5mm);
            \fill[black] ([shift={(-7mm,-7mm)}]table.south west) rectangle ++(5mm,5mm);
            \fill[black] ([shift={(7mm,-7mm)}]table.south east) rectangle ++(-5mm,5mm);
        \end{tikzpicture}
    \end{center}
    \renewcommand{\arraystretch}{1.0}
}

% Chamada do gabarito
\Gabarito

\newpage
\end{document}